\section{Preconditioned (quadratic-metric) SAGA (hypothesis B1)}
\label{sec:B1}

Throughout, $F(\theta)=\frac1n\sum_{i=1}^n f_i(\theta)$ on $\R^d$, each $f_i$ is convex and $R^2$-smooth,
$F$ is $\mu$-strongly convex with minimizer $\etastar$, $\kappa=R^2/\mu$, and $P\in\R^{d\times d}$ is a
\emph{fixed} symmetric positive definite (SPD) matrix; $P^{1/2}$ denotes its SPD square root and
$P^{-1/2}=(P^{1/2})^{-1}$. Whenever a Hessian is written we assume the corresponding function is twice
continuously differentiable (this is the case for logistic regression); every statement that only involves
gradients holds without this assumption (Remark~\ref{rem:B1-C1}). We write $\Ei$ for the expectation with
respect to the index $i(t)$ drawn at iteration $t$, conditionally on the past.

The content of this note is the following chain of facts.
(1)~SAGA run in the metric of $P$ (``$P$-SAGA'') is \emph{exactly} Euclidean SAGA on a linearly
transformed problem (Proposition~\ref{prop:B1-transfer}), so Proposition~5.9 of \cite{bach2024learning}
transfers verbatim, with $(R^2,\mu)$ replaced by the relative constants $(R_P^2,\mu_P)$ of
\eqref{eq:B1-RP-muP} (Corollary~\ref{cor:B1-rate}). (2)~Redoing the proof in the $P$-metric shows which of
the six proof inequalities change (Section~\ref{sec:B1-proofP}); this defines the six tightness ratios to be
instrumented in the $P$-geometry (Definition~\ref{def:B1-ratios}). (3)~On $\ell_2$-regularized logistic
regression, \emph{no} quadratic metric improves the certified global condition number, $\kappaP\ge\kappa$ for
every $P$ (Proposition~\ref{prop:B1-nofreelunch}): the benefit of preconditioning is entirely a
trajectory-dependent effect, captured by the effective constants of the B2 note. For the statistical
preconditioner $P=\hat\Sigma+\lambda I$ the trajectory-dependent condition number is at most
$(d_\infty(\lambda)/4+1)/m(\theta)$, where $d_\infty(\lambda)$ ($n$ times the largest ridge leverage score, at
most $C\,d_{\mathrm{eff}}(\lambda)$ under a coherence assumption) is the only trace of the spectrum of
$\hat\Sigma$ and $m(\theta)$ is the \emph{smallest} per-sample logistic curvature averaged on the segment
$[\etastar,\theta]$ (Corollary~\ref{cor:B1-kappa-traj}). Because of the factor $1/m(\theta)$ this rigorous
bound need not improve on the Euclidean one, and on the instances of our experiments it does not: it exceeds
$\kappa$, i.e.\ it predicts no gain. The observed gain is accounted for only by the local constants at the
optimum (Lemma~\ref{lem:B1-asymptotic}), whose use along the whole run is heuristic
(Section~\ref{sec:B1-predictions}). (4)~Cost and the case of a non-constant metric
(Section~\ref{sec:B1-cost}), and a prediction table for the experiments (Section~\ref{sec:B1-predictions}).

\subsection{The \texorpdfstring{$P$}{P}-geometry and the algorithm}
\label{sec:B1-algo}

\begin{definition}[$P$-norm and its dual norm]\label{def:B1-norms}
For $v,g\in\R^d$ let $\norm{v}_P^2:=v^\top P v$ (norm on parameters) and
$\norm{g}_{P^{-1}}^2:=g^\top P^{-1} g$ (norm on gradients and memory vectors).
\end{definition}

\begin{lemma}[elementary facts on the $P$-geometry]\label{lem:B1-geometry}
Let $P$ be SPD. Then:
\begin{enumerate}[label=(\roman*)]
\item (generalized Cauchy--Schwarz) $|\inner{g}{v}|\le \norm{g}_{P^{-1}}\norm{v}_P$ for all $g,v\in\R^d$,
with equality if and only if $g$ and $Pv$ are collinear; consequently
$\norm{g}_{P^{-1}}=\sup_{\norm{v}_P\le 1}\inner{g}{v}$, i.e. $\norm{\cdot}_{P^{-1}}$ is the dual norm of
$\norm{\cdot}_P$;
\item $\norm{P^{1/2}v}=\norm{v}_P$, $\norm{P^{-1/2}g}=\norm{g}_{P^{-1}}$ and $\norm{P^{-1}g}_P=\norm{g}_{P^{-1}}$;
\item (square-norm inequality) $\norm{a+b}_{P^{-1}}^2\le 2\norm{a}_{P^{-1}}^2+2\norm{b}_{P^{-1}}^2$,
with equality iff $a=b$;
\item (variance $\le$ second moment) for a random vector $Z$ with $\E\norm{Z}^2<\infty$,
$\E\norm{Z-\E Z}_{P^{-1}}^2=\E\norm{Z}_{P^{-1}}^2-\norm{\E Z}_{P^{-1}}^2\le \E\norm{Z}_{P^{-1}}^2$;
\item $\lmin(P)\norm{v}^2\le\norm{v}_P^2\le\lmax(P)\norm{v}^2$ and
$\lmax(P)^{-1}\norm{g}^2\le\norm{g}_{P^{-1}}^2\le\lmin(P)^{-1}\norm{g}^2$.
\end{enumerate}
\end{lemma}
\begin{proof}
(i) $\inner{g}{v}=\inner{P^{-1/2}g}{P^{1/2}v}$, and the Euclidean Cauchy--Schwarz inequality gives
$|\inner{g}{v}|\le\norm{P^{-1/2}g}\,\norm{P^{1/2}v}=\norm{g}_{P^{-1}}\norm{v}_P$, with equality iff
$P^{-1/2}g$ and $P^{1/2}v$ are collinear, i.e. iff $g$ and $Pv$ are; the choice
$v=P^{-1}g/\norm{g}_{P^{-1}}$ (for $g\neq0$) has $\norm{v}_P=1$ and attains the supremum.
(ii) $\norm{P^{1/2}v}^2=v^\top Pv$, $\norm{P^{-1/2}g}^2=g^\top P^{-1}g$ and
$(P^{-1}g)^\top P(P^{-1}g)=g^\top P^{-1}g$.
(iii) $\norm{a+b}_{P^{-1}}^2=\norm{a}_{P^{-1}}^2+2a^\top P^{-1}b+\norm{b}_{P^{-1}}^2$ and
$2a^\top P^{-1}b\le\norm{a}_{P^{-1}}^2+\norm{b}_{P^{-1}}^2$ because
$\norm{a-b}_{P^{-1}}^2\ge0$, with equality iff $a=b$.
(iv) Expand $\E\norm{Z-\E Z}_{P^{-1}}^2=\E\norm{Z}_{P^{-1}}^2-2(\E Z)^\top P^{-1}\E Z+\norm{\E Z}_{P^{-1}}^2$.
(v) These are the extremal values of the Rayleigh quotients of $P$ and of $P^{-1}$.
\end{proof}

\begin{definition}[$P$-SAGA: SAGA in the metric of $P$]\label{def:B1-PSAGA}
Fix $\gamma>0$, $\theta_0\in\R^d$ and memory vectors $z_i^{(0)}\in\R^d$, $i=1,\dots,n$. At iteration
$t\ge1$, draw $i(t)$ uniformly in $\{1,\dots,n\}$, independently of the past, and set
\begin{align}
g_t&:=f'_{i(t)}(\theta_{t-1})+\frac1n\sum_{i=1}^n z_i^{(t-1)}-z_{i(t)}^{(t-1)},\label{eq:B1-gt}\\
\theta_t&:=\argmin_{\theta\in\R^d}\Big\{\inner{g_t}{\theta}
+\frac1{2\gamma}(\theta-\theta_{t-1})^\top P(\theta-\theta_{t-1})\Big\},\label{eq:B1-argmin}\\
z_{i(t)}^{(t)}&:=f'_{i(t)}(\theta_{t-1}),\qquad z_j^{(t)}:=z_j^{(t-1)}\quad(j\neq i(t)).\label{eq:B1-memory}
\end{align}
The memory update \eqref{eq:B1-memory} is the one of SAGA \cite{defazio2014saga}; only the step
\eqref{eq:B1-argmin} depends on $P$.
\end{definition}

\begin{lemma}[closed form, unbiasedness, Bregman form]\label{lem:B1-step}
The minimization \eqref{eq:B1-argmin} has the unique solution
\begin{equation}\label{eq:B1-step}
\theta_t=\theta_{t-1}-\gamma P^{-1}g_t ,
\end{equation}
which for $P=I$ is the SAGA iteration of \cite[Sec.~5.4.4]{bach2024learning}. Moreover
$\Ei[g_t]=\Fp(\theta_{t-1})$. Finally, \eqref{eq:B1-argmin} is the Bregman step
$\argmin_\theta\{\gamma\inner{g_t}{\theta}+D_h(\theta,\theta_{t-1})\}$ for the quadratic reference function
$h=\frac12\norm{\cdot}_P^2$, whose Bregman divergence is $D_h(\theta,\theta')=\frac12\norm{\theta-\theta'}_P^2$;
i.e. $P$-SAGA is Algorithm~1 (Bregman-SAGA) of \cite{dragomir2021fast} with a quadratic $h$.
\end{lemma}
\begin{proof}
The objective in \eqref{eq:B1-argmin} is a strictly convex quadratic (its Hessian is $P/\gamma\succ0$), so
it has a unique minimizer, characterized by a vanishing gradient: $g_t+\gamma^{-1}P(\theta-\theta_{t-1})=0$,
i.e. $\theta=\theta_{t-1}-\gamma P^{-1}g_t$. Since $i(t)$ is uniform and independent of the past,
$\Ei[f'_{i(t)}(\theta_{t-1})]=\frac1n\sum_i\fip(\theta_{t-1})=\Fp(\theta_{t-1})$ and
$\Ei[z_{i(t)}^{(t-1)}]=\frac1n\sum_i z_i^{(t-1)}$, so the two memory terms of \eqref{eq:B1-gt} cancel in
expectation. For $h=\frac12\norm{\cdot}_P^2$ one has $\nabla h(\theta')=P\theta'$ and
$D_h(\theta,\theta')=h(\theta)-h(\theta')-\inner{\nabla h(\theta')}{\theta-\theta'}
=\frac12\norm{\theta-\theta'}_P^2$; multiplying the objective of \eqref{eq:B1-argmin} by $\gamma$ gives the
Bregman form.
\end{proof}

\subsection{Change of variables: \texorpdfstring{$P$}{P}-SAGA is SAGA on a transformed problem}
\label{sec:B1-transfer}

\begin{proposition}[$P$-SAGA as SAGA on a transformed problem]\label{prop:B1-transfer}
Let $u:=P^{1/2}\theta$, $\tilde f_i(u):=f_i(P^{-1/2}u)$ and $\tilde F:=\frac1n\sum_i\tilde f_i=F\circ P^{-1/2}$.
Then:
\begin{enumerate}[label=(\roman*)]
\item $\nabla\tilde f_i(u)=P^{-1/2}\fip(P^{-1/2}u)$ and, when $f_i$ is $C^2$,
$\nabla^2\tilde f_i(u)=P^{-1/2}f_i''(P^{-1/2}u)P^{-1/2}$; the same formulas hold for $(\tilde F,F)$.
\item $\tilde F$ is strongly convex and $\tilde u_*:=P^{1/2}\etastar$ is its unique minimizer.
\item Let $(\theta_t,z^{(t)})_{t\ge0}$ be generated by $P$-SAGA (Definition~\ref{def:B1-PSAGA}) with step
$\gamma$ and index sequence $(i(t))_{t\ge1}$, and let $u_t:=P^{1/2}\theta_t$,
$\tilde z_i^{(t)}:=P^{-1/2}z_i^{(t)}$. Then $(u_t,\tilde z^{(t)})_{t\ge0}$ is generated by Euclidean SAGA
($P=I$) on $(\tilde f_i)_{i\le n}$ with the \emph{same} step $\gamma$ and the \emph{same} index sequence,
started from $(u_0,\tilde z^{(0)})$.
\item For all $t$ and $i$,
$\norm{u_t-\tilde u_*}^2=\norm{\theta_t-\etastar}_P^2$ and
$\norm{\tilde z_i^{(t)}-\nabla\tilde f_i(\tilde u_*)}^2=\norm{z_i^{(t)}-\fip(\etastar)}_{P^{-1}}^2$;
in particular $z_i^{(0)}=\fip(\theta_0)$ for all $i$ iff $\tilde z_i^{(0)}=\nabla\tilde f_i(u_0)$ for all $i$.
\end{enumerate}
\end{proposition}
\begin{proof}
(i) $\tilde f_i=f_i\circ A$ with the linear map $A=P^{-1/2}$ (symmetric), so by the chain rule
$\nabla\tilde f_i(u)=A^\top\fip(Au)=P^{-1/2}\fip(P^{-1/2}u)$ and
$\nabla^2\tilde f_i(u)=A^\top f_i''(Au)A$. Averaging over $i$ gives the statements for $\tilde F$.

(ii) By (i), $\nabla\tilde F(\tilde u_*)=P^{-1/2}\Fp(\etastar)=0$. For $u,u'\in\R^d$ write
$\theta=P^{-1/2}u$, $\theta'=P^{-1/2}u'$; by (i) and $\mu$-strong convexity of $F$,
\begin{align*}
\tilde F(u')-\tilde F(u)-\inner{\nabla\tilde F(u)}{u'-u}
&=F(\theta')-F(\theta)-\inner{\Fp(\theta)}{\theta'-\theta}\\
&\ge\frac\mu2\norm{\theta'-\theta}^2\ge\frac{\mu}{2\lmax(P)}\norm{u'-u}^2 ,
\end{align*}
where we used $\inner{P^{-1/2}\Fp(\theta)}{u'-u}=\inner{\Fp(\theta)}{\theta'-\theta}$ and
Lemma~\ref{lem:B1-geometry}(v), which gives
$\norm{\theta'-\theta}^2=\norm{P^{-1/2}(u'-u)}^2\ge\norm{u'-u}^2/\lmax(P)$.
Hence $\tilde F$ is $(\mu/\lmax(P))$-strongly convex and its stationary point $\tilde u_*$ is its unique
minimizer.

(iii) By induction on $t$; the initialization is the definition of $(u_0,\tilde z^{(0)})$. Let
$\tilde g_t:=\nabla\tilde f_{i(t)}(u_{t-1})+\frac1n\sum_i\tilde z_i^{(t-1)}-\tilde z_{i(t)}^{(t-1)}$ be the
Euclidean SAGA direction built from $(u_{t-1},\tilde z^{(t-1)})$ and the index $i(t)$. By (i) and the
induction hypothesis,
$\tilde g_t=P^{-1/2}f'_{i(t)}(\theta_{t-1})+P^{-1/2}\big(\frac1n\sum_i z_i^{(t-1)}-z_{i(t)}^{(t-1)}\big)
=P^{-1/2}g_t$. Hence, by \eqref{eq:B1-step},
$u_t=P^{1/2}\theta_t=P^{1/2}\theta_{t-1}-\gamma P^{1/2}P^{-1}g_t=u_{t-1}-\gamma P^{-1/2}g_t
=u_{t-1}-\gamma\tilde g_t$, which is the Euclidean SAGA step. For the memory,
$\tilde z_{i(t)}^{(t)}=P^{-1/2}f'_{i(t)}(\theta_{t-1})=\nabla\tilde f_{i(t)}(u_{t-1})$ and
$\tilde z_j^{(t)}=\tilde z_j^{(t-1)}$ for $j\neq i(t)$: this is the SAGA memory update.

(iv) By Lemma~\ref{lem:B1-geometry}(ii),
$\norm{u_t-\tilde u_*}^2=\norm{P^{1/2}(\theta_t-\etastar)}^2=\norm{\theta_t-\etastar}_P^2$ and
$\norm{\tilde z_i^{(t)}-\nabla\tilde f_i(\tilde u_*)}^2=\norm{P^{-1/2}(z_i^{(t)}-\fip(\etastar))}^2
=\norm{z_i^{(t)}-\fip(\etastar)}_{P^{-1}}^2$, using (i) at $\tilde u_*$. The last claim is (i) at $u_0$,
$P^{-1/2}$ being invertible.
\end{proof}

\begin{lemma}[relative smoothness and strong convexity constants]\label{lem:B1-constants}
\begin{enumerate}[label=(\roman*)]
\item $\tilde f_i$ is $L$-smooth (i.e. $\nabla\tilde f_i$ is $L$-Lipschitz) if and only if
\begin{equation}\label{eq:B1-rel-smooth}
\norm{\fip(\theta)-\fip(\theta')}_{P^{-1}}\le L\,\norm{\theta-\theta'}_P\qquad\forall\theta,\theta'\in\R^d ;
\end{equation}
when $f_i$ is $C^2$ and convex, this is equivalent to $\sup_{\theta}\lmax(P^{-1/2}f_i''(\theta)P^{-1/2})\le L$.
\item $\tilde F$ is $m$-strongly convex if and only if
$F(\theta')\ge F(\theta)+\inner{\Fp(\theta)}{\theta'-\theta}+\frac m2\norm{\theta'-\theta}_P^2$ for all
$\theta,\theta'$; when $F$ is $C^2$, this is equivalent to $\inf_\theta\lmin(P^{-1/2}F''(\theta)P^{-1/2})\ge m$.
\item If each $f_i$ is $R^2$-smooth and $F$ is $\mu$-strongly convex (Euclidean constants), then
\eqref{eq:B1-rel-smooth} holds with $L=R^2/\lmin(P)$ and (ii) holds with $m=\mu/\lmax(P)$. Conversely, if
\eqref{eq:B1-rel-smooth} holds with $L=R_P^2$ and (ii) with $m=\mu_P$, then each $f_i$ is
$\lmax(P)R_P^2$-smooth and $F$ is $\lmin(P)\mu_P$-strongly convex.
\end{enumerate}
\end{lemma}
\begin{proof}
(i) With $\theta=P^{-1/2}u$, $\theta'=P^{-1/2}u'$, Proposition~\ref{prop:B1-transfer}(i) and
Lemma~\ref{lem:B1-geometry}(ii) give
\[
\norm{\nabla\tilde f_i(u)-\nabla\tilde f_i(u')}=\norm{P^{-1/2}(\fip(\theta)-\fip(\theta'))}
=\norm{\fip(\theta)-\fip(\theta')}_{P^{-1}},\qquad \norm{u-u'}=\norm{\theta-\theta'}_P ;
\]
as $u$ ranges over $\R^d$ so does $\theta$. For a $C^2$ convex function, $L$-smoothness is equivalent to all Hessians being
$\preceq LI$ \cite[Def.~5.3 and Sec.~5.2]{bach2024learning}; apply this to $\tilde f_i$, whose Hessian at
$u$ is $P^{-1/2}f_i''(P^{-1/2}u)P^{-1/2}$. (ii) Same substitution in the definition of strong convexity,
using $\inner{\nabla\tilde F(u)}{u'-u}=\inner{\Fp(\theta)}{\theta'-\theta}$ and
$\norm{u'-u}^2=\norm{\theta'-\theta}_P^2$; for $C^2$ functions $m$-strong convexity is equivalent to all
Hessians being $\succeq mI$ \cite[Sec.~5.2]{bach2024learning}.
(iii) By Lemma~\ref{lem:B1-geometry}(v),
\begin{align*}
\norm{\fip(\theta)-\fip(\theta')}_{P^{-1}}&\le\lmin(P)^{-1/2}\norm{\fip(\theta)-\fip(\theta')}
\le\lmin(P)^{-1/2}R^2\norm{\theta-\theta'}\\
&\le\lmin(P)^{-1}R^2\norm{\theta-\theta'}_P ,
\end{align*}
and $\frac\mu2\norm{\theta'-\theta}^2\ge\frac{\mu}{2\lmax(P)}\norm{\theta'-\theta}_P^2$. The converse uses the
reverse inequalities of Lemma~\ref{lem:B1-geometry}(v) in the same way.
\end{proof}

\begin{definition}[relative constants and relative condition number]\label{def:B1-kappaP}
We set
\begin{equation}\label{eq:B1-RP-muP}
\begin{gathered}
R_P^2:=\max_{1\le i\le n}\ \sup_{\theta\in\R^d}\lmax\big(P^{-1/2}f_i''(\theta)P^{-1/2}\big),\\
\mu_P:=\inf_{\theta\in\R^d}\lmin\big(P^{-1/2}F''(\theta)P^{-1/2}\big),\qquad
\kappaP:=\frac{R_P^2}{\mu_P} .
\end{gathered}
\end{equation}
By Lemma~\ref{lem:B1-constants}, every $\tilde f_i$ of Proposition~\ref{prop:B1-transfer} is $R_P^2$-smooth
and $\tilde F$ is $\mu_P$-strongly convex; in the terminology of \cite[Def.~1]{hendrikx2020statistically} and
\cite{lu2018relatively}, $f_i$ is $R_P^2$-smooth and $F$ is $\mu_P$-strongly convex \emph{relative to} the reference
function $h=\frac12\norm{\cdot}_P^2$ (whose Hessian is $P$). For $P=I$ these are the usual constants, and
$R^2,\mu,\kappa$ always denote the Euclidean ($P=I$) constants.
\end{definition}

\begin{lemma}[comparison with the Euclidean constants]\label{lem:B1-compare}
Let $\operatorname{cond}(P):=\lmax(P)/\lmin(P)$ be the condition number of $P$. Then
\begin{gather}
\mu_P\le\sup_\theta\lmax\big(P^{-1/2}F''(\theta)P^{-1/2}\big)\le R_P^2,\qquad
R_P^2\le\frac{R^2}{\lmin(P)},\qquad \mu_P\ge\frac{\mu}{\lmax(P)},\nonumber\\
\frac{\kappa}{\operatorname{cond}(P)}\le\kappaP\le\kappa\,\operatorname{cond}(P).\label{eq:B1-kappaP-vs-kappa}
\end{gather}
In particular $\kappaP\ge1$, $\kappa_{cP}=\kappaP$ for every $c>0$, and a badly chosen metric can degrade
the condition number by at most $\operatorname{cond}(P)$.
\end{lemma}
\begin{proof}
$F''=\frac1n\sum_i f_i''$ and $\lmax$ is subadditive on symmetric matrices, so
$\lmax(P^{-1/2}F''(\theta)P^{-1/2})\le\frac1n\sum_i\lmax(P^{-1/2}f_i''(\theta)P^{-1/2})\le R_P^2$ for every
$\theta$, while $\lmin\le\lmax$ at every $\theta$ gives the first inequality. The two middle inequalities
are Lemma~\ref{lem:B1-constants}(iii) combined with the Hessian characterizations; dividing them gives the
upper bound on $\kappaP$, and the converse part of Lemma~\ref{lem:B1-constants}(iii) gives
$\kappa\le\kappaP\operatorname{cond}(P)$. Finally $R_{cP}^2=R_P^2/c$ and $\mu_{cP}=\mu_P/c$ by
\eqref{eq:B1-RP-muP}.
\end{proof}

\begin{remark}[scaling of $P$ and of the step]\label{rem:B1-scaling}
$(cP)$-SAGA with step $c\gamma$ produces exactly the same iterates as $P$-SAGA with step $\gamma$, since
$c\gamma(cP)^{-1}=\gamma P^{-1}$; consistently, the normalized step $\gamma=1/(4R_P^2)$ of
Corollary~\ref{cor:B1-rate} is multiplied by $c$ when $P$ is. Only $\kappaP$ is scale free; the values of
$R_P^2$ and $\mu_P$ quoted below for $P=\hat\Sigma+\lambda I$ depend on the normalization
$\hat\Sigma=\frac1nX^\top X$ (with $P'=nP=X^\top X+n\lambda I$ one would get $R_{P'}^2=R_P^2/n$).
\end{remark}

\begin{remark}[$C^1$ functions]\label{rem:B1-C1}
If the $f_i$ are only $C^1$, define $R_P^2$ as the smallest constant $L$ in \eqref{eq:B1-rel-smooth}
(over all $i$) and $\mu_P$ as the largest $m$ in Lemma~\ref{lem:B1-constants}(ii). Everything below that
does not display a Hessian remains valid, because the transfer argument and the proof of Proposition~5.9
only use the gradient forms (co-coercivity, Proposition~5.4 of \cite{bach2024learning}, is proved from
$L$-smoothness and convexity alone).
\end{remark}

\begin{lemma}[relative constants of logistic regression]\label{lem:B1-logistic}
Let $f_i(\theta)=\ell(y_ix_i^\top\theta)+\frac\lambda2\norm{\theta}^2$ with $\ell(u)=\log(1+e^{-u})$,
$y_i\in\{-1,1\}$, $x_i\in\R^d$, $\lambda>0$, and let $\sigma(u)=(1+e^{-u})^{-1}$. Then
$\ell'(u)=-\sigma(-u)$ and $\ell''(u)=\sigma'(u)=\sigma(u)\sigma(-u)=\big(2+e^{u}+e^{-u}\big)^{-1}\in(0,\tfrac14]$,
which equals $\frac14$ at $u=0$ and is strictly decreasing in $|u|$. Writing
$\sigma_i'(\theta):=\sigma'(y_ix_i^\top\theta)$,
\begin{equation}\label{eq:B1-logistic-hessians}
\fip(\theta)=\ell'(y_ix_i^\top\theta)\,y_ix_i+\lambda\theta,\qquad
f_i''(\theta)=\sigma_i'(\theta)\,x_ix_i^\top+\lambda I,\qquad
F''(\theta)=\frac1n\sum_{i=1}^n\sigma_i'(\theta)\,x_ix_i^\top+\lambda I .
\end{equation}
For every SPD $P$,
\begin{equation}\label{eq:B1-logistic-RP}
\max\Big\{\frac14\max_i\norm{x_i}_{P^{-1}}^2,\ \frac{\max_i\norm{x_i}^2/4+\lambda}{\lmax(P)}\Big\}\ \le\ R_P^2\ \le\
\frac14\max_i\norm{x_i}_{P^{-1}}^2+\frac{\lambda}{\lmin(P)} ,
\end{equation}
and more precisely $R_P^2\ge\frac14\norm{x_i}_{P^{-1}}^2+\lambda\,\frac{x_i^\top P^{-2}x_i}{x_i^\top P^{-1}x_i}$
for every $i$ with $x_i\neq0$;
\begin{equation}\label{eq:B1-logistic-muP}
\mu_P=\lambda\,\lmin(P^{-1})=\frac{\lambda}{\lmax(P)} .
\end{equation}
In particular, for $P=I$: $R^2:=\frac14\max_i\norm{x_i}^2+\lambda$ is a valid smoothness constant of
every $f_i$, $\mu=\lambda$, and $\kappa=R^2/\mu=\max_i\norm{x_i}^2/(4\lambda)+1$; these are the Euclidean
constants used in the whole project.
\end{lemma}
\begin{proof}
$\ell'(u)=-e^{-u}/(1+e^{-u})=-\sigma(-u)$ and $\ell''(u)=\sigma'(u)=\sigma(u)(1-\sigma(u))=\sigma(u)\sigma(-u)
=1/(2+e^u+e^{-u})$, which is maximal at $u=0$ (value $\frac14$) and decreasing in $|u|$ since
$e^u+e^{-u}=2\cosh u$ is increasing in $|u|$. The gradient and Hessians follow from the chain rule and
$y_i^2=1$. Upper bound in \eqref{eq:B1-logistic-RP}: by Weyl's inequality
$\lmax(A+B)\le\lmax(A)+\lmax(B)$ with $A=\sigma_i'(\theta)P^{-1/2}x_ix_i^\top P^{-1/2}$ (rank one, with
nonzero eigenvalue $\sigma_i'(\theta)\norm{P^{-1/2}x_i}^2\le\frac14\norm{x_i}_{P^{-1}}^2$) and
$B=\lambda P^{-1}$. Lower bounds: at $\theta=0$, $\sigma_i'(0)=\frac14$; for $x_i\neq0$ the Rayleigh
quotient of $P^{-1/2}f_i''(0)P^{-1/2}$ along $v=P^{-1/2}x_i$ equals
$\frac14(x_i^\top P^{-1}x_i)^2/\norm{P^{-1/2}x_i}^2+\lambda\norm{P^{-1}x_i}^2/\norm{P^{-1/2}x_i}^2
=\frac14\norm{x_i}_{P^{-1}}^2+\lambda\,(x_i^\top P^{-2}x_i)/(x_i^\top P^{-1}x_i)$, which is the ``more
precise'' bound; it is $\ge\frac14\norm{x_i}_{P^{-1}}^2$, and also $\ge(\norm{x_i}^2/4+\lambda)/\lmax(P)$
because, with $y:=P^{-1/2}x_i$, the fraction $(x_i^\top P^{-2}x_i)/(x_i^\top P^{-1}x_i)=y^\top P^{-1}y/\norm y^2$
is a Rayleigh quotient of $P^{-1}$, hence $\ge\lmin(P^{-1})=1/\lmax(P)$, while
$\norm{x_i}_{P^{-1}}^2\ge\norm{x_i}^2/\lmax(P)$ by Lemma~\ref{lem:B1-geometry}(v). If $x_i=0$ then
$f_i''\equiv\lambda I$ and $\sup_\theta\lmax(P^{-1/2}f_i''(\theta)P^{-1/2})=\lambda/\lmin(P)\ge(\norm{x_i}^2/4+\lambda)/\lmax(P)$,
the first lower bound being trivial. Taking the maximum over $i$ gives the two lower bounds in
\eqref{eq:B1-logistic-RP}.
For \eqref{eq:B1-logistic-muP}: $F''(\theta)\succeq\lambda I$ for all $\theta$ gives
$P^{-1/2}F''(\theta)P^{-1/2}\succeq\lambda P^{-1}\succeq\lambda\lmin(P^{-1})I$, hence
$\mu_P\ge\lambda\lmin(P^{-1})$. Conversely, let $I_0=\{i:x_i\neq0\}$ and pick $v$ with $v^\top x_i\neq0$
for all $i\in I_0$ (the excluded set is a finite union of hyperplanes). For $s\to\infty$,
$|y_ix_i^\top(sv)|\to\infty$ hence $\sigma_i'(sv)\to0$ for $i\in I_0$, while the terms $i\notin I_0$
vanish identically; thus $F''(sv)\to\lambda I$ and, by continuity of $\lmin$,
$\lmin(P^{-1/2}F''(sv)P^{-1/2})\to\lambda\lmin(P^{-1})$, so the infimum in \eqref{eq:B1-RP-muP} is
$\le\lambda\lmin(P^{-1})$. The $P=I$ statements are the special case $\lmin(P)=\lmax(P)=1$.
\end{proof}

\subsection{Convergence of \texorpdfstring{$P$}{P}-SAGA: Proposition~5.9 in the \texorpdfstring{$P$}{P}-metric}
\label{sec:B1-rate}

\begin{corollary}[Proposition~5.9 transferred to $P$-SAGA]\label{cor:B1-rate}
Let $P$ be SPD and let $R_P^2,\mu_P$ be the constants of \eqref{eq:B1-RP-muP} (or any constants such that
each $\tilde f_i$ of Proposition~\ref{prop:B1-transfer} is $R_P^2$-smooth and $\tilde F$ is $\mu_P$-strongly
convex). Run $P$-SAGA (Definition~\ref{def:B1-PSAGA}) with $\gamma=1/(4R_P^2)$ and define
\begin{gather}
V_t^P:=\norm{\theta_t-\etastar}_P^2+3\gamma^2\sum_{i=1}^n\norm{z_i^{(t)}-\fip(\etastar)}_{P^{-1}}^2,\nonumber\\
\rho_P:=1-\min\Big\{\frac1{3n},\frac{3\mu_P}{16R_P^2}\Big\}=1-\min\Big\{\frac1{3n},\frac{3}{16\kappaP}\Big\}.\label{eq:B1-VP}
\end{gather}
Then for every $t\ge1$,
\begin{equation}\label{eq:B1-onestep}
\Ei\big[V_t^P\big]\le\rho_P\,V_{t-1}^P ,
\end{equation}
and if $z_i^{(0)}=\fip(\theta_0)$ for all $i$,
\begin{equation}\label{eq:B1-rate}
\E\norm{\theta_t-\etastar}_P^2\le\rho_P^{\,t}\Big(1+\frac n4\Big)\norm{\theta_0-\etastar}_P^2 .
\end{equation}
Consequently
\[
\E\norm{\theta_t-\etastar}^2\le\operatorname{cond}(P)\,\rho_P^{\,t}\Big(1+\frac n4\Big)\norm{\theta_0-\etastar}^2
\]
(the condition number of $P$ enters the constant, not the rate), and
$t\ge\max\{3n,\frac{16}{3}\kappaP\}\log\frac{2n}{\varepsilon}$ iterations guarantee
$\E\norm{\theta_t-\etastar}_P^2\le\varepsilon\norm{\theta_0-\etastar}_P^2$, i.e. a running time of order
$c_P\,(3n+\frac{16}{3}\kappaP)\log\frac{2n}{\varepsilon}$, where $c_P$ is the cost of one iteration
(Section~\ref{sec:B1-cost}).
\end{corollary}

\begin{proof}[Proof by transfer]
By Proposition~\ref{prop:B1-transfer}(iii), $(u_t,\tilde z^{(t)})=(P^{1/2}\theta_t,P^{-1/2}z^{(t)})$ is the
Euclidean SAGA sequence on $(\tilde f_i)$ with step $\gamma$; the $\tilde f_i$ are convex and
$R_P^2$-smooth and $\tilde F$ is $\mu_P$-strongly convex with minimizer $\tilde u_*$
(Lemma~\ref{lem:B1-constants}, Proposition~\ref{prop:B1-transfer}(ii)). The proof of Proposition~5.9 in
\cite{bach2024learning} (transcribed in the B3 note) establishes, for Euclidean SAGA with $\gamma=1/(4L)$ on
any family of convex $L$-smooth functions whose average is $m$-strongly convex, the one-step contraction
$\Ei[\tilde V_t]\le(1-\min\{\frac1{3n},\frac{3m}{16L}\})\tilde V_{t-1}$ of the Lyapunov function
$\tilde V_t=\norm{u_t-\tilde u_*}^2+3\gamma^2\sum_i\norm{\tilde z_i^{(t)}-\nabla\tilde f_i(\tilde u_*)}^2$,
and then the bound $\E\norm{u_t-\tilde u_*}^2\le(1-\min\{\frac1{3n},\frac{3m}{16L}\})^t(1+\frac n4)
\norm{u_0-\tilde u_*}^2$ when $\tilde z_i^{(0)}=\nabla\tilde f_i(u_0)$. By
Proposition~\ref{prop:B1-transfer}(iv), $\tilde V_t=V_t^P$ and
$\norm{u_t-\tilde u_*}^2=\norm{\theta_t-\etastar}_P^2$ identically, and the initialization conditions
correspond. The Euclidean consequence follows from Lemma~\ref{lem:B1-geometry}(v) applied on both sides;
the iteration count follows from $\rho_P^t\le\exp(-t\min\{\frac1{3n},\frac3{16\kappaP}\})$ and
$1+\frac n4\le 2n$.
\end{proof}

\subsubsection{The same proof written in the \texorpdfstring{$P$}{P}-metric: which inequalities change}
\label{sec:B1-proofP}

To see exactly where $P$ enters, we rewrite the proof of Proposition~5.9 in the original coordinates;
every inequality is labelled, and these labels are the ones instrumented by the experiments
(Definition~\ref{def:B1-ratios}). Fix $t\ge1$ and write $\delta:=\theta_{t-1}-\etastar$,
$\omega_t:=g_t$, $\bar z:=\frac1n\sum_i z_i^{(t-1)}$ and, for each $i$,
\[
a_{t,i}:=\fip(\theta_{t-1})-\fip(\etastar),\qquad
b_{t,i}:=\fip(\etastar)-z_i^{(t-1)}+\bar z,\qquad\text{so that}\quad \omega_t=a_{t,i(t)}+b_{t,i(t)} .
\]

\emph{Step 1.} Expanding the square in the $P$-inner product and using \eqref{eq:B1-step},
\begin{align}
\norm{\theta_t-\etastar}_P^2=\norm{\delta-\gamma P^{-1}\omega_t}_P^2
&=\norm{\delta}_P^2-2\gamma\,\delta^\top PP^{-1}\omega_t+\gamma^2\omega_t^\top P^{-1}PP^{-1}\omega_t\nonumber\\
&=\norm{\delta}_P^2-2\gamma\inner{\delta}{\omega_t}+\gamma^2\norm{\omega_t}_{P^{-1}}^2 .\label{eq:B1-expand}
\end{align}
The pairing $\inner{\delta}{\omega_t}$ is the same number for every $P$; taking $\Ei$ and using
Lemma~\ref{lem:B1-step} (unbiasedness, also metric independent),
$\Ei\norm{\theta_t-\etastar}_P^2=\norm{\delta}_P^2-2\gamma\inner{\delta}{\Fp(\theta_{t-1})}
+\gamma^2\Ei\norm{\omega_t}_{P^{-1}}^2$. The three inequalities of Step~1 are
\begin{align}
\text{(I1$_P$, square norm)}\quad&
\Ei\norm{\omega_t}_{P^{-1}}^2\le 2\,\Ei\norm{a_{t,i(t)}}_{P^{-1}}^2+2\,\Ei\norm{b_{t,i(t)}}_{P^{-1}}^2,
\label{ineq:P-sq}\\
\text{(I2$_P$, co-coercivity)}\quad&
\Ei\norm{a_{t,i(t)}}_{P^{-1}}^2=\frac1n\sum_{i=1}^n\norm{\fip(\theta_{t-1})-\fip(\etastar)}_{P^{-1}}^2\nonumber\\
&\qquad\le R_P^2\,\inner{\Fp(\theta_{t-1})}{\theta_{t-1}-\etastar},
\label{ineq:P-coco}\\
\text{(I3$_P$, variance $\le$ 2nd moment)}\quad&
\Ei\norm{b_{t,i(t)}}_{P^{-1}}^2\le\frac1n\sum_{i=1}^n\norm{\fip(\etastar)-z_i^{(t-1)}}_{P^{-1}}^2 .
\label{ineq:P-var}
\end{align}
\eqref{ineq:P-sq} is Lemma~\ref{lem:B1-geometry}(iii) applied for each value of $i(t)$. For
\eqref{ineq:P-coco}: each $\tilde f_i$ is convex and $R_P^2$-smooth, so co-coercivity (Proposition~5.4 of
\cite{bach2024learning}) gives
$\norm{\nabla\tilde f_i(u)-\nabla\tilde f_i(u')}^2\le R_P^2\inner{\nabla\tilde f_i(u)-\nabla\tilde f_i(u')}{u-u'}$;
with $u=P^{1/2}\theta_{t-1}$, $u'=P^{1/2}\etastar$ and Proposition~\ref{prop:B1-transfer}(i), this reads
\begin{equation}\label{eq:B1-coco-i}
\norm{\fip(\theta_{t-1})-\fip(\etastar)}_{P^{-1}}^2\le R_P^2\inner{\fip(\theta_{t-1})-\fip(\etastar)}{\theta_{t-1}-\etastar}
\qquad(i=1,\dots,n),
\end{equation}
because $\inner{P^{-1/2}(\fip(\theta)-\fip(\etastar))}{P^{1/2}(\theta-\etastar)}
=\inner{\fip(\theta)-\fip(\etastar)}{\theta-\etastar}$; averaging over $i$ and using
$\sum_i\fip(\etastar)=n\Fp(\etastar)=0$ gives \eqref{ineq:P-coco}, which is (5.31) of the book in the
$P$-metric. For \eqref{ineq:P-var}: with $Z:=f'_{i(t)}(\etastar)-z_{i(t)}^{(t-1)}$ one has $\Ei Z=-\bar z$
(again $\sum_i\fip(\etastar)=0$), so $Z-\Ei Z=b_{t,i(t)}$ and Lemma~\ref{lem:B1-geometry}(iv) applies.
Combining, Step~1 yields
\begin{equation}\label{eq:B1-step1}
\Ei\norm{\theta_t-\etastar}_P^2\le\norm{\delta}_P^2-2\gamma(1-\gamma R_P^2)\inner{\Fp(\theta_{t-1})}{\delta}
+\frac{2\gamma^2}{n}\sum_{i=1}^n\norm{\fip(\etastar)-z_i^{(t-1)}}_{P^{-1}}^2 .
\end{equation}

\emph{Step 2.} Only the $i(t)$-th memory vector changes, so, exactly,
\begin{align*}
\sum_{i=1}^n\norm{\fip(\etastar)-z_i^{(t)}}_{P^{-1}}^2&=\sum_{i=1}^n\norm{\fip(\etastar)-z_i^{(t-1)}}_{P^{-1}}^2
-\norm{f'_{i(t)}(\etastar)-z_{i(t)}^{(t-1)}}_{P^{-1}}^2\\
&\quad+\norm{f'_{i(t)}(\etastar)-f'_{i(t)}(\theta_{t-1})}_{P^{-1}}^2 ;
\end{align*}
taking $\Ei$ (each index has probability $1/n$) and using \eqref{ineq:P-coco} once more,
\begin{align}
\Ei\sum_{i=1}^n\norm{\fip(\etastar)-z_i^{(t)}}_{P^{-1}}^2
&=\Big(1-\frac1n\Big)\sum_{i=1}^n\norm{\fip(\etastar)-z_i^{(t-1)}}_{P^{-1}}^2
+\frac1n\sum_{i=1}^n\norm{\fip(\etastar)-\fip(\theta_{t-1})}_{P^{-1}}^2\nonumber\\
&\le\Big(1-\frac1n\Big)\sum_{i=1}^n\norm{\fip(\etastar)-z_i^{(t-1)}}_{P^{-1}}^2
+R_P^2\inner{\Fp(\theta_{t-1})}{\delta}.\label{eq:B1-step2}
\end{align}

\emph{Combination.} Multiply \eqref{eq:B1-step2} by $\Delta=3\gamma^2$ and add to \eqref{eq:B1-step1}.
With $\gamma=1/(4R_P^2)$: $1-\gamma R_P^2-R_P^2\Delta/(2\gamma)=1-\frac14-\frac38=\frac38$ and
$\frac{2\gamma^2}{n\Delta}+1-\frac1n=1-\frac1{3n}$, so
\begin{equation}\label{eq:B1-combined}
\Ei[V_t^P]\le\norm{\delta}_P^2-\frac{3\gamma}{4}\inner{\Fp(\theta_{t-1})}{\delta}
+\Big(1-\frac1{3n}\Big)\Delta\sum_{i=1}^n\norm{\fip(\etastar)-z_i^{(t-1)}}_{P^{-1}}^2 .
\end{equation}
The last two inequalities are
\begin{align}
\text{(I4$_P$, strong convexity)}\quad&
\inner{\Fp(\theta_{t-1})}{\theta_{t-1}-\etastar}\ \ge\ \mu_P\,\norm{\theta_{t-1}-\etastar}_P^2,
\label{ineq:P-sc}\\
\text{(I5$_P$, min)}\quad&
\Big(1-\frac{3\mu_P}{16R_P^2}\Big)\norm{\delta}_P^2
+\Big(1-\frac1{3n}\Big)\Delta\sum_{i=1}^n\norm{\fip(\etastar)-z_i^{(t-1)}}_{P^{-1}}^2\ \le\ \rho_P\,V_{t-1}^P .
\label{ineq:P-min}
\end{align}
For \eqref{ineq:P-sc}: $\tilde F$ is $\mu_P$-strongly convex, so adding the inequalities of
Lemma~\ref{lem:B1-constants}(ii) written at $(\theta,\theta')=(\theta_{t-1},\etastar)$ and at
$(\etastar,\theta_{t-1})$, and using $\Fp(\etastar)=0$, gives
$\inner{\Fp(\theta_{t-1})}{\theta_{t-1}-\etastar}\ge\mu_P\norm{\theta_{t-1}-\etastar}_P^2$; inserting it in
\eqref{eq:B1-combined} (the coefficient $-\frac{3\gamma}4$ is negative, so the inequality direction is
preserved) produces the left-hand side of \eqref{ineq:P-min}, since $\frac{3\gamma}{4}\mu_P=\frac{3\mu_P}{16R_P^2}$.
\eqref{ineq:P-min} is the scalar inequality $(1-\alpha)x+(1-\beta)y\le(1-\min\{\alpha,\beta\})(x+y)$ for
$x,y\ge0$, with $\alpha=\frac{3\mu_P}{16R_P^2}$, $\beta=\frac1{3n}$. This proves \eqref{eq:B1-onestep}.
Iterating with the tower property, $\E V_t^P\le\rho_P^tV_0^P$, and
\begin{equation}\label{ineq:P-init}
\text{(I6$_P$, initialization)}\qquad
\sum_{i=1}^n\norm{\fip(\etastar)-\fip(\theta_0)}_{P^{-1}}^2\le n\,R_P^4\,\norm{\theta_0-\etastar}_P^2
\end{equation}
(each term is bounded by the square of \eqref{eq:B1-rel-smooth}) gives, when $z_i^{(0)}=\fip(\theta_0)$,
\[
V_0^P\le(1+3\gamma^2nR_P^4)\norm{\theta_0-\etastar}_P^2=\Big(1+\frac{3n}{16}\Big)\norm{\theta_0-\etastar}_P^2
\le\Big(1+\frac n4\Big)\norm{\theta_0-\etastar}_P^2 ,
\]
which is \eqref{eq:B1-rate} because
$\norm{\theta_t-\etastar}_P^2\le V_t^P$. This completes the second proof of
Corollary~\ref{cor:B1-rate}. \qed

\begin{remark}[comparison with Bregman-SAGA]\label{rem:B1-dragomir}
Corollary~1 of \cite{dragomir2021fast} states, for Bregman-SAGA with a quadratic reference function
($\nabla^2h$ constant, which is exactly $P$-SAGA by Lemma~\ref{lem:B1-step}) and step $1/(8L_{f/h})$,
the contraction $\E[\psi_t]\le(1-\min\{\frac{1}{8\kappa_{f/h}},\frac1{2n}\})^t\psi_0$ of their potential
$\psi_t=\frac1{\eta}D_h(x^\star,x_t)+\frac n2\cdot\frac1n\sum_iD_{f_i}(\phi_i^t,x^\star)$, where
$L_{f/h}$ is the relative smoothness of every $f_i$ (their Assumption~3) and
$\kappa_{f/h}=L_{f/h}/\mu_{f/h}$ the relative condition number: these are our $R_P^2$ and $\kappaP$. The
rate has the same structure $1-\min\{c_1/n,\,c_2/\kappaP\}$ as \eqref{eq:B1-VP} (with $(c_1,c_2)=(\frac12,\frac18)$
instead of $(\frac13,\frac3{16})$, the difference coming from the step and from a potential expressed with
Bregman distances of the $f_i$ instead of squared norms of memory vectors).
\end{remark}

\subsubsection{Metric dependence of each step, and the six tightness ratios in the \texorpdfstring{$P$}{P}-geometry}
\label{sec:B1-ratios}

Table~\ref{tab:B1-metric} summarizes where $P$ enters. Three kinds of objects appear.
(a) Quantities that are the same numbers for every $P$, given the current state $(\theta_{t-1},z^{(t-1)})$:
the pairing $\inner{\theta_{t-1}-\etastar}{\omega_t}$ and its conditional expectation
$\inner{\theta_{t-1}-\etastar}{\Fp(\theta_{t-1})}$, the unbiasedness, the memory bookkeeping identity of
Step~2 (an identity, valid for any norm), the scalar inequality (I5), and the algebra of the combination
(the constants $\frac38$, $\frac2{3n}$, $\frac3{16}$ are unchanged because $\gamma$ and $\Delta$ are
normalized by $R_P^2$). (b) Inequalities whose \emph{form} is metric independent but which are instantiated
in the $P^{-1}$-inner product: (I1) and (I3), both valid in every Hilbert norm (Lemma~\ref{lem:B1-geometry}).
(c) Inequalities whose norm \emph{and} constant change: (I2) (left side in $\norm{\cdot}_{P^{-1}}$, right
side metric free, constant $R_P^2$), (I4) (left side metric free, right side in $\norm{\cdot}_P$, constant
$\mu_P$), (I6) (both norms and $R_P^4$). Of course the iterates themselves depend on $P$, so the numerical
values of all ratios depend on $P$; the point of the classification is that the \emph{same six ratios} can
be computed, by Proposition~\ref{prop:B1-transfer}(iv), either in $\theta$-coordinates with $P$-norms (one
linear solve with $P$ per vector), or in $u$-coordinates with Euclidean norms on the transformed problem
(whose regularizer is the non-isotropic quadratic $\frac\lambda2u^\top P^{-1}u$ for logistic regression,
so the B3 code must accept a general quadratic regularizer if the second route is taken).

\begin{table}[ht]
\centering
\small
\begin{tabular}{@{}p{2.0cm}p{2.5cm}p{4.8cm}p{5.2cm}@{}}
\toprule
Step & Inequality & What changes with $P$ & Metric-free part \\
\midrule
(I1) \eqref{ineq:P-sq} & square norm & norm $\norm{\cdot}_{P^{-1}}$ only & form of the inequality \\
(I2) \eqref{ineq:P-coco} & co-coercivity & norm $\norm{\cdot}_{P^{-1}}$, constant $R_P^2$ &
  pairing $\inner{\Fp(\theta_{t-1})}{\theta_{t-1}-\etastar}$ \\
(I3) \eqref{ineq:P-var} & variance $\le$ 2nd moment & norm $\norm{\cdot}_{P^{-1}}$ only & form of the inequality \\
(I4) \eqref{ineq:P-sc} & strong convexity & norm $\norm{\cdot}_P$, constant $\mu_P$ &
  pairing $\inner{\Fp(\theta_{t-1})}{\theta_{t-1}-\etastar}$ \\
(I5) \eqref{ineq:P-min} & min & only through the value of $3\mu_P/(16R_P^2)$ & scalar inequality \\
(I6) \eqref{ineq:P-init} & initialization & both norms, constant $R_P^4$ & --- \\
exact steps & \eqref{eq:B1-expand}, Step~2 identity, combination & norms only & unbiasedness, pairing, algebra \\
\bottomrule
\end{tabular}
\caption{Metric dependence of the steps of the proof of Proposition~5.9 (Section~\ref{sec:B1-proofP}).}
\label{tab:B1-metric}
\end{table}

\begin{definition}[tightness ratios in the $P$-geometry]\label{def:B1-ratios}
The effective (trajectory-dependent) constants of the B2 note, written in the $P$-metric, are, for
$\theta\neq\etastar$,
\begin{equation}\label{eq:B1-eff-constants}
\muEff^P(\theta):=\frac{\inner{\Fp(\theta)}{\theta-\etastar}}{\norm{\theta-\etastar}_P^2},
\qquad
\LEff^P(\theta):=\frac{\frac1n\sum_i\norm{\fip(\theta)-\fip(\etastar)}_{P^{-1}}^2}
{\inner{\Fp(\theta)}{\theta-\etastar}} .
\end{equation}
\emph{Convention on indices.} As in the B3 note, the ratios of iteration $t$ are evaluated at the state
$(\theta_{t-1},z^{(t-1)})$ seen before $i(t)$ is drawn, so they involve $\muEff^P(\theta_{t-1})$ and
$\LEff^P(\theta_{t-1})$; the B2 note abbreviates $\muEff(t):=\muEff(\theta_t)$, which for $P=I$ is the
value entering the ratios of iteration $t+1$. To avoid any off-by-one error in code shared by the three
notes, we never write $\muEff^P(t)$ but always display the point, $\muEff^P(\theta_{t-1})$.
With the notation of Section~\ref{sec:B1-proofP} ($\delta=\theta_{t-1}-\etastar$, $a_{t,i}$, $b_{t,i}$,
$\bar z$), $\alpha=\frac{3\mu_P}{16R_P^2}$, $\beta=\frac1{3n}$ and $\Delta=3\gamma^2$, define, for
$\theta_{t-1}\neq\etastar$,
\begin{align*}
r^P_{\mathrm{sq}}(t)&:=\frac{\frac1n\sum_i\norm{a_{t,i}+b_{t,i}}_{P^{-1}}^2}
{\frac2n\sum_i\big(\norm{a_{t,i}}_{P^{-1}}^2+\norm{b_{t,i}}_{P^{-1}}^2\big)},\\
r^P_{\mathrm{coco}}(t)&:=\frac{\frac1n\sum_i\norm{\fip(\theta_{t-1})-\fip(\etastar)}_{P^{-1}}^2}
{R_P^2\,\inner{\Fp(\theta_{t-1})}{\delta}}=\frac{\LEff^P(\theta_{t-1})}{R_P^2},\\
r^P_{\mathrm{var}}(t)&:=\frac{\frac1n\sum_i\norm{b_{t,i}}_{P^{-1}}^2}
{\frac1n\sum_i\norm{\fip(\etastar)-z_i^{(t-1)}}_{P^{-1}}^2}
=1-\frac{\norm{\bar z}_{P^{-1}}^2}{\frac1n\sum_i\norm{\fip(\etastar)-z_i^{(t-1)}}_{P^{-1}}^2},\\
r^P_{\mathrm{sc}}(t)&:=\frac{\mu_P\norm{\delta}_P^2}{\inner{\Fp(\theta_{t-1})}{\delta}}
=\frac{\mu_P}{\muEff^P(\theta_{t-1})},\\
r^P_{\mathrm{min}}(t)&:=\frac{(1-\alpha)\norm{\delta}_P^2+(1-\beta)\Delta\sum_i\norm{\fip(\etastar)-z_i^{(t-1)}}_{P^{-1}}^2}
{\rho_P\,V_{t-1}^P},
\end{align*}
with the convention $r^P_{\mathrm{var}}(t):=1$ when the memory is exact ($z_i^{(t-1)}=\fip(\etastar)$ for
all $i$, the only case in which its denominator vanishes; \eqref{ineq:P-var} is then the equality
$0\le0$), and, for $\theta_0\neq\etastar$,
\begin{equation}\label{eq:B1-rinit-identity}
r^P_{\mathrm{init}}:=\frac{\sum_i\norm{\fip(\etastar)-\fip(\theta_0)}_{P^{-1}}^2}
{nR_P^4\norm{\theta_0-\etastar}_P^2}
=\frac{\LEff^P(\theta_0)\,\muEff^P(\theta_0)}{R_P^4}
=r^P_{\mathrm{coco}}(1)\,\frac{\muEff^P(\theta_0)}{R_P^2} .
\end{equation}
Each ratio is the smaller side of one of \eqref{ineq:P-sq}--\eqref{ineq:P-init} divided by the larger side,
the expectations over $i(t)$ being exact averages over $i$. The identities in \eqref{eq:B1-rinit-identity}
are immediate from \eqref{eq:B1-eff-constants}, the pairing $\inner{\Fp(\theta_0)}{\theta_0-\etastar}$
cancelling in the product $\LEff^P(\theta_0)\muEff^P(\theta_0)$: $r^P_{\mathrm{init}}$ is a derived ratio,
not an independent one. The one-step ratio $r^P_{\mathrm{tot}}(t):=\Ei[V_t^P]/(\rho_PV_{t-1}^P)$, defined
when $V_{t-1}^P>0$, is also computable exactly (average of $V_t^P$ over the $n$ possible draws). For $P=I$
all of them reduce to the ratios of the B3 note.
\end{definition}

\begin{lemma}[ranges and degenerate cases]\label{lem:B1-ranges}
Let $\theta_{t-1}\neq\etastar$. Then
\begin{equation}\label{eq:B1-eff-sandwich}
0<\mu_P\le\muEff^P(\theta_{t-1})\le\LEff^P(\theta_{t-1})\le R_P^2 ,
\end{equation}
all the denominators in Definition~\ref{def:B1-ratios} are positive (that of $r^P_{\mathrm{var}}$ when the
memory is not exact), and
$r^P_{\mathrm{sc}}(t),r^P_{\mathrm{coco}}(t),r^P_{\mathrm{sq}}(t),r^P_{\mathrm{min}}(t)\in(0,1]$,
$r^P_{\mathrm{var}}(t)\in[0,1]$; if $\theta_0\neq\etastar$ then $r^P_{\mathrm{init}}\in(0,1]$. At
$\theta_{t-1}=\etastar$ the ratios $r^P_{\mathrm{sc}}$ and $r^P_{\mathrm{coco}}$ are of the form $0/0$ (both
numerator and denominator vanish since $\Fp(\etastar)=0$), as are $r^P_{\mathrm{sq}}$ when in addition all
memory errors $z_i^{(t-1)}-\fip(\etastar)$ coincide, and $r^P_{\mathrm{min}}$, $r^P_{\mathrm{tot}}$ when in
addition the memory is exact; they are left undefined there, as in the B3 note.
\end{lemma}
\begin{proof}
By Proposition~\ref{prop:B1-transfer}(i),(iv), every quantity above is the corresponding Euclidean quantity
of the B3 note computed on the transformed problem $(\tilde f_i)$ at the state $(u_{t-1},\tilde z^{(t-1)})$,
for which the hypotheses of the B3 note hold (convex $R_P^2$-smooth $\tilde f_i$, $\mu_P$-strongly convex
$\tilde F$, $u_{t-1}\neq\tilde u_*$); the inequalities of the statement are Lemma~\ref{lem:B3-ranges} of the B3
note (ranges of the tightness ratios), transferred. For completeness, directly: \eqref{ineq:P-sc} gives
$\inner{\Fp(\theta_{t-1})}{\delta}\ge\mu_P\norm{\delta}_P^2>0$, i.e. $\muEff^P(\theta_{t-1})\ge\mu_P>0$, so
$\Fp(\theta_{t-1})\neq0$; by convexity of the squared norm (Jensen) and Lemma~\ref{lem:B1-geometry}(i),
$\frac1n\sum_i\norm{\fip(\theta_{t-1})-\fip(\etastar)}_{P^{-1}}^2\ge\norm{\Fp(\theta_{t-1})}_{P^{-1}}^2
\ge\inner{\Fp(\theta_{t-1})}{\delta}^2/\norm{\delta}_P^2>0$, which is $\LEff^P\ge\muEff^P$ together with the
positivity of the numerator of $r^P_{\mathrm{coco}}$; $\LEff^P\le R_P^2$ is \eqref{ineq:P-coco}. The
numerator of $r^P_{\mathrm{sq}}$ is $\ge\norm{\frac1n\sum_i(a_{t,i}+b_{t,i})}_{P^{-1}}^2
=\norm{\Fp(\theta_{t-1})}_{P^{-1}}^2>0$ (Jensen again, and $\sum_ib_{t,i}=0$ because
$\sum_i\fip(\etastar)=0$); the numerator of $r^P_{\mathrm{min}}$ is $\ge(1-\alpha)\norm{\delta}_P^2>0$ as
$\alpha\le\frac3{16}$, and its denominator is $\rho_PV_{t-1}^P\ge\rho_P\norm{\delta}_P^2>0$. The upper
bounds are the inequalities \eqref{ineq:P-sq}--\eqref{ineq:P-min} themselves, and $r^P_{\mathrm{var}}\ge0$
is clear. For $r^P_{\mathrm{init}}$, apply \eqref{eq:B1-eff-sandwich} at $\theta_0$ in
\eqref{eq:B1-rinit-identity}.
\end{proof}

\begin{remark}[effective one-step contraction]\label{rem:B1-effective-rate}
Keep $\gamma=1/(4R_P^2)$ and $\Delta=3\gamma^2$ but, in Steps~1--2, replace the two uses of
\eqref{ineq:P-coco} and the use of \eqref{ineq:P-sc} by the identities
\[
\frac1n\sum_i\norm{\fip(\theta_{t-1})-\fip(\etastar)}_{P^{-1}}^2=\LEff^P(\theta_{t-1})\inner{\Fp(\theta_{t-1})}{\delta},
\qquad
\inner{\Fp(\theta_{t-1})}{\delta}=\muEff^P(\theta_{t-1})\norm{\delta}_P^2 .
\]
The coefficient of $\inner{\Fp(\theta_{t-1})}{\delta}$ in the combination becomes (writing $\LEff^P$ for
$\LEff^P(\theta_{t-1})$)
\[
-2\gamma\big(1-\gamma\LEff^P-\LEff^P\Delta/(2\gamma)\big)=-2\gamma\big(1-\tfrac52\gamma\LEff^P\big)
=-2\gamma\big(1-\tfrac58r^P_{\mathrm{coco}}(t)\big),
\]
whence, with only \eqref{ineq:P-sq} and \eqref{ineq:P-var} still used,
\begin{equation}\label{eq:B1-effective-rate}
\Ei[V_t^P]\le\Big(1-\Big(1-\frac58r^P_{\mathrm{coco}}(t)\Big)\frac{\muEff^P(\theta_{t-1})}{2R_P^2}\Big)
\norm{\theta_{t-1}-\etastar}_P^2
+\Big(1-\frac1{3n}\Big)\Delta\sum_{i=1}^n\norm{\fip(\etastar)-z_i^{(t-1)}}_{P^{-1}}^2 .
\end{equation}
Since $r^P_{\mathrm{coco}}\in(0,1]$ (Lemma~\ref{lem:B1-ranges}), the contraction factor of the parameter
part lies between $1-\frac{\muEff^P(\theta_{t-1})}{2R_P^2}$ and $1-\frac{3\muEff^P(\theta_{t-1})}{16R_P^2}$:
the rate ``seen by the proof'' is governed by the effective relative condition number
$R_P^2/\muEff^P(\theta_{t-1})$ (the B2 quantity in the
$P$-metric), the looseness of co-coercivity affecting only the constant, between $\frac3{16}$ and
$\frac12$. This is the quantity whose behaviour under the four metrics is predicted in
Section~\ref{sec:B1-predictions}.
\end{remark}

\subsection{Choice of \texorpdfstring{$P$}{P} on \texorpdfstring{$\ell_2$}{l2}-regularized logistic regression}
\label{sec:B1-choice}

\paragraph{Setting and notation.}
$X\in\R^{n\times d}$ has rows $x_i^\top$, $\hat\Sigma:=\frac1nX^\top X$ has eigenvalues
$\sigma_1\ge\dots\ge\sigma_d\ge0$ and eigenvectors $v_1,\dots,v_d$, and $f_i$ is as in
Lemma~\ref{lem:B1-logistic}. We use the ridge leverage scores and the effective dimension
(\cite[Def.~1]{elalaoui2015fast} with the linear kernel $K=XX^\top$; \cite[Sec.~4.1]{bach2013sharp}):
\begin{align}
\mathrm{lev}_i(\lambda)&:=x_i^\top(X^\top X+n\lambda I)^{-1}x_i=\big[X(X^\top X+n\lambda I)^{-1}X^\top\big]_{ii},\label{eq:B1-leverage-def}\\
d_{\mathrm{eff}}(\lambda)&:=\tr\big(\hat\Sigma(\hat\Sigma+\lambda I)^{-1}\big)=\sum_{k=1}^d\frac{\sigma_k}{\sigma_k+\lambda},\qquad
d_{\infty}(\lambda):=n\max_i\mathrm{lev}_i(\lambda),\nonumber
\end{align}
$d_\infty$ being the ``maximal marginal degrees of freedom''
$n\,\|\diag(K(K+n\lambda I)^{-1})\|_\infty$ of \cite[Sec.~4.1, eq.~(5)]{bach2013sharp}.

\emph{Notation shared with the other notes.} The segment average $\bar\sigma_i'(\theta)$ of
Proposition~\ref{prop:B1-traj-sc} below is the mean-value weight $w_i(\theta)$ of the B2 note, and
$w_i:=\sigma_i'(\etastar)$ (paragraph~(d)) is its value $w_i(\etastar)$ at the optimum, denoted $p_i$ in the B2
note; our $m(\theta)=\min_i\bar\sigma_i'(\theta)$ is at least the $m(\theta)$ of the B2 note (the minimum of
$\sigma'$ over the segments, attained at an endpoint), so every bound below stated with $m(\theta)$ also holds
with the latter. The coherence constants $C$ (Assumption~\ref{ass:B1-coherence}) and $C_w$ (paragraph~(d)) are
unrelated to the cross term $C$ of the B3 note, and the relative Hessian $M_P$ of
Section~\ref{sec:B1-predictions} is unrelated to the memory error $M$ of the B3 note and to the matrix
$M_\ast$ of the B2 note. Inside this note, $m$ also denotes a generic strong-convexity constant
(Lemma~\ref{lem:B1-constants}) and the inner-loop length in \cite{moritz2016linearly,gower2016stochastic}.

\begin{lemma}[leverage scores]\label{lem:B1-leverage}
\begin{enumerate}[label=(\roman*)]
\item $\norm{x_i}^2_{(\hat\Sigma+\lambda I)^{-1}}=n\,\mathrm{lev}_i(\lambda)$ for every $i$.
\item $\sum_{i=1}^n\mathrm{lev}_i(\lambda)=d_{\mathrm{eff}}(\lambda)$; hence
$d_{\mathrm{eff}}(\lambda)\le d_\infty(\lambda)$.
\item $0\le\mathrm{lev}_i(\lambda)<1$ for every $i$, hence $d_\infty(\lambda)<n$; for every $i$ the map
$\lambda\mapsto\mathrm{lev}_i(\lambda)$ is nonincreasing on $(0,\infty)$, hence so are $d_\infty$ and
$d_{\mathrm{eff}}$; moreover
$d_{\mathrm{eff}}(\lambda)\le\min\{\operatorname{rank}(X),\,\tr(\hat\Sigma)/\lambda\}$.
\item $X(X^\top X+n\lambda I)^{-1}X^\top=K(K+n\lambda I)^{-1}$ with $K=XX^\top$, so $\mathrm{lev}_i(\lambda)$ is
the $\lambda$-ridge leverage score of \cite[Def.~1]{elalaoui2015fast}.
\end{enumerate}
\end{lemma}
\begin{proof}
(i) $(\hat\Sigma+\lambda I)^{-1}=\big(\frac1n(X^\top X+n\lambda I)\big)^{-1}=n(X^\top X+n\lambda I)^{-1}$.
(ii) By cyclicity of the trace,
$\sum_i\mathrm{lev}_i=\tr\big(X(X^\top X+n\lambda I)^{-1}X^\top\big)=\tr\big((X^\top X+n\lambda I)^{-1}X^\top X\big)
=\tr\big((\hat\Sigma+\lambda I)^{-1}\hat\Sigma\big)$, which equals $\sum_k\sigma_k/(\sigma_k+\lambda)$ in
the eigenbasis of $\hat\Sigma$; the inequality is $\max_i\mathrm{lev}_i\ge\frac1n\sum_i\mathrm{lev}_i$.
(iii) Let $X=\sum_{k=1}^{r}s_ku_kw_k^\top$ be a thin SVD ($r=\operatorname{rank}X$, $s_k>0$,
$u_k\in\R^n$ orthonormal). Then $X(X^\top X+n\lambda I)^{-1}X^\top=\sum_k\frac{s_k^2}{s_k^2+n\lambda}u_ku_k^\top$,
whose $i$-th diagonal entry is $\mathrm{lev}_i=\sum_k\frac{s_k^2}{s_k^2+n\lambda}u_{ki}^2$. If
$\sum_ku_{ki}^2=0$ then $\mathrm{lev}_i=0$; otherwise $\frac{s_k^2}{s_k^2+n\lambda}<1$ gives
$\mathrm{lev}_i<\sum_ku_{ki}^2\le1$, the last inequality because the $u_k$ are orthonormal columns of an
$n\times r$ matrix. Each coefficient $s_k^2/(s_k^2+n\lambda)$ is decreasing in $\lambda$ and $u_{ki}^2\ge0$,
so $\mathrm{lev}_i(\lambda)$ is nonincreasing in $\lambda$, and so are $d_\infty=n\max_i\mathrm{lev}_i$ and
$d_{\mathrm{eff}}=\sum_i\mathrm{lev}_i$. The bound on $d_{\mathrm{eff}}$ follows from
$\sigma_k/(\sigma_k+\lambda)\le\min\{\ind_{\sigma_k>0},\sigma_k/\lambda\}$.
(iv) Push-through identity: $X(X^\top X+n\lambda I)=(XX^\top+n\lambda I)X$, hence
$(XX^\top+n\lambda I)^{-1}X=X(X^\top X+n\lambda I)^{-1}$; multiply by $X^\top$ on the right.
\end{proof}

\begin{assumption}[coherence of the leverage scores]\label{ass:B1-coherence}
There is $C\ge1$ such that the ridge leverage scores \eqref{eq:B1-leverage-def} satisfy
$\max_i\mathrm{lev}_i(\lambda)\le C\,d_{\mathrm{eff}}(\lambda)/n$, i.e.
$d_\infty(\lambda)\le C\,d_{\mathrm{eff}}(\lambda)$: no data point has a leverage score larger than $C$
times the average leverage score.
\end{assumption}

\paragraph{(a) $P=I$ (SAGA).}
By Lemma~\ref{lem:B1-logistic}, $R^2=\frac14\max_i\norm{x_i}^2+\lambda$, $\mu=\lambda$ and
\begin{equation}\label{eq:B1-kappa-I}
\kappa=\frac{\max_i\norm{x_i}^2}{4\lambda}+1\ \ge\ \frac{\tr(\hat\Sigma)}{4\lambda}+1
=\frac{\sum_{k=1}^d\sigma_k}{4\lambda}+1 ,
\end{equation}
since $\max_i\norm{x_i}^2\ge\frac1n\sum_i\norm{x_i}^2=\tr\hat\Sigma$. Thus $\kappa$ contains the whole
spectrum of $\hat\Sigma$ through its trace (not only $\sigma_1$, which is what the condition number
$L/\mu$ of $F$ alone would contain, since $L\le\sigma_1/4+\lambda$), divided by $\lambda$.

\paragraph{(b) $P=P_{\mathrm J}:=\diag(\hat\Sigma)+\lambda I$ (feature scaling, Jacobi preconditioning).}
Write $\hat\Sigma_{jj}=\frac1n\sum_ix_{ij}^2$ for the diagonal entries of $\hat\Sigma$, assume no feature is
identically zero (otherwise drop it), and let
$C_{\mathrm{corr}}:=\diag(\hat\Sigma)^{-1/2}\,\hat\Sigma\,\diag(\hat\Sigma)^{-1/2}$ be the correlation matrix
(unit diagonal, so $\lmin(C_{\mathrm{corr}})\le1\le\lmax(C_{\mathrm{corr}})\le d$).
\begin{proposition}[constants of the Jacobi preconditioner]\label{prop:B1-diag}
\begin{enumerate}[label=(\roman*)]
\item $\norm{x_i}_{P_{\mathrm J}^{-1}}^2=\sum_{j=1}^d\frac{x_{ij}^2}{\hat\Sigma_{jj}+\lambda}$ and
\[
\frac1n\sum_i\norm{x_i}_{P_{\mathrm J}^{-1}}^2=\sum_{j=1}^d\frac{\hat\Sigma_{jj}}{\hat\Sigma_{jj}+\lambda}=:d^{\diag}_{\mathrm{eff}}(\lambda)
\ \ge\ d_{\mathrm{eff}}(\lambda).
\]
\item $R_{P_{\mathrm J}}^2\le\frac14\max_i\norm{x_i}_{P_{\mathrm J}^{-1}}^2+\frac{\lambda}{\min_j\hat\Sigma_{jj}+\lambda}$ and
$\mu_{P_{\mathrm J}}=\frac{\lambda}{\max_j\hat\Sigma_{jj}+\lambda}$.
\item $\lmin(C_{\mathrm{corr}})\,P_{\mathrm J}\preceq\hat\Sigma+\lambda I\preceq\lmax(C_{\mathrm{corr}})\,P_{\mathrm J}$;
consequently $\lmin(C_{\mathrm{corr}})\,n\,\mathrm{lev}_i(\lambda)\le\norm{x_i}_{P_{\mathrm J}^{-1}}^2
\le\lmax(C_{\mathrm{corr}})\,n\,\mathrm{lev}_i(\lambda)$ for all $i$.
\end{enumerate}
\end{proposition}
\begin{proof}
(i) $P_{\mathrm J}$ is diagonal, so $x_i^\top P_{\mathrm J}^{-1}x_i=\sum_jx_{ij}^2/(\hat\Sigma_{jj}+\lambda)$; averaging over $i$ and
using $\frac1n\sum_ix_{ij}^2=\hat\Sigma_{jj}$ gives $d^{\diag}_{\mathrm{eff}}$. For the inequality,
$g(s)=s/(s+\lambda)$ is concave on $[0,\infty)$, the vector $(\hat\Sigma_{jj})_j$ of diagonal entries of
$\hat\Sigma$ is majorized by the vector of its eigenvalues (Schur--Horn theorem), and
$x\mapsto\sum_jg(x_j)$ is Schur-concave when $g$ is concave \cite{marshall2011inequalities}; hence
$\sum_jg(\hat\Sigma_{jj})\ge\sum_kg(\sigma_k)$.
(ii) Lemma~\ref{lem:B1-logistic} with $\lmin(P_{\mathrm J})=\min_j\hat\Sigma_{jj}+\lambda$ and
$\lmax(P_{\mathrm J})=\max_j\hat\Sigma_{jj}+\lambda$.
(iii) $\hat\Sigma=\diag(\hat\Sigma)^{1/2}C_{\mathrm{corr}}\diag(\hat\Sigma)^{1/2}$, so
$\lmin(C_{\mathrm{corr}})\diag(\hat\Sigma)\preceq\hat\Sigma\preceq\lmax(C_{\mathrm{corr}})\diag(\hat\Sigma)$; adding $\lambda I$ and
using $\lmin(C_{\mathrm{corr}})\le1\le\lmax(C_{\mathrm{corr}})$ gives the sandwich. If $\lmin(C_{\mathrm{corr}})>0$,
inverting it (the map $A\mapsto A^{-1}$ is order reversing on SPD matrices) gives
\[
\lmax(C_{\mathrm{corr}})^{-1}P_{\mathrm J}^{-1}\preceq(\hat\Sigma+\lambda I)^{-1}\preceq\lmin(C_{\mathrm{corr}})^{-1}P_{\mathrm J}^{-1},
\]
which with Lemma~\ref{lem:B1-leverage}(i) yields the bounds on $\norm{x_i}_{P_{\mathrm J}^{-1}}^2$; if
$\lmin(C_{\mathrm{corr}})=0$ (singular $\hat\Sigma$), the lower bound is trivial and the upper bound follows by
inverting the right-hand inequality alone.
\end{proof}
Thus the Jacobi metric removes the \emph{scale} anisotropy of the features (up to the $\lambda$ terms, it
is invariant under a rescaling of each feature) but keeps their \emph{correlation} anisotropy: every
constant of the statistical preconditioner $\hat\Sigma+\lambda I$ derived below transfers to $P_{\mathrm J}$ at the
price of the factor $\operatorname{cond}(C_{\mathrm{corr}})=\lmax(C_{\mathrm{corr}})/\lmin(C_{\mathrm{corr}})$
(the argument of Lemma~\ref{lem:B1-compare}, which applies verbatim to any pair of metrics, applied to
$(P_{\mathrm J},\hat\Sigma+\lambda I)$ through (iii)), and $P_{\mathrm J}=\hat\Sigma+\lambda I$
exactly when the features are uncorrelated ($C_{\mathrm{corr}}=I$).

\paragraph{(c) $P=\hat\Sigma+\lambda I$ (statistical preconditioning).}
Here $\lmax(P)=\sigma_1+\lambda$, $\lmin(P)=\sigma_d+\lambda$ and, by Lemma~\ref{lem:B1-leverage}(i),
$\norm{x_i}_{P^{-1}}^2=n\,\mathrm{lev}_i(\lambda)$.

\begin{proposition}[global constants of the statistical preconditioner]\label{prop:B1-stat-global}
For $P=\hat\Sigma+\lambda I$,
\begin{equation}\label{eq:B1-stat-RP}
\frac{d_\infty(\lambda)}{4}\ \le\ R_P^2\ \le\ \frac{d_\infty(\lambda)}{4}+\frac{\lambda}{\sigma_d+\lambda}
\ \le\ \frac{d_\infty(\lambda)}4+1\ \le\ \frac{C\,d_{\mathrm{eff}}(\lambda)}{4}+1,
\end{equation}
the last inequality under Assumption~\ref{ass:B1-coherence}, and
\begin{equation}\label{eq:B1-stat-muP}
\mu_P=\frac{\lambda}{\sigma_1+\lambda},\qquad
\kappaP=\frac{R_P^2}{\mu_P}\in\Big[\frac{d_\infty(\lambda)}{4}\cdot\frac{\sigma_1+\lambda}{\lambda},\
\Big(\frac{d_\infty(\lambda)}{4}+\frac{\lambda}{\sigma_d+\lambda}\Big)\frac{\sigma_1+\lambda}{\lambda}\Big].
\end{equation}
\end{proposition}
\begin{proof}
Insert $\max_i\norm{x_i}_{P^{-1}}^2=n\max_i\mathrm{lev}_i(\lambda)=d_\infty(\lambda)$ and $\lmin(P)=\sigma_d+\lambda$
in \eqref{eq:B1-logistic-RP}, and $\lmax(P)=\sigma_1+\lambda$ in \eqref{eq:B1-logistic-muP};
$\lambda/(\sigma_d+\lambda)\le1$ and Assumption~\ref{ass:B1-coherence} give the last two bounds.
\end{proof}

The smoothness side is now controlled by leverage scores, i.e. by an effective dimension, with no trace of
the spectrum of $\hat\Sigma$; but the \emph{global} strong convexity side has become
$\lambda/(\sigma_1+\lambda)$, and the following statement shows that this is not an artefact of the proof:
with global constants, no quadratic metric can beat $P=I$ on this problem.

\begin{proposition}[no free lunch for global constants]\label{prop:B1-nofreelunch}
For $\ell_2$-regularized logistic regression and \emph{every} SPD matrix $P$,
\begin{equation}\label{ineq:P-nofreelunch}
\kappaP\ \ge\ \frac{\max_i\norm{x_i}^2}{4\lambda}+1=\kappa ,
\end{equation}
with $\kappa$ the Euclidean condition number \eqref{eq:B1-kappa-I}; consequently $\rho_P\ge\rho_I$ for every
$P$. In particular, for $P=\hat\Sigma+\lambda I$, $\kappa\le\kappaP\le\kappa\,\frac{\sigma_1+\lambda}{\sigma_d+\lambda}$.
\end{proposition}
\begin{proof}
By the second lower bound in \eqref{eq:B1-logistic-RP},
$R_P^2\ge(\max_i\norm{x_i}^2/4+\lambda)/\lmax(P)=R^2/\lmax(P)$, while $\mu_P=\lambda/\lmax(P)$ by
\eqref{eq:B1-logistic-muP}; divide. Since $\kappaP\mapsto\rho_P=1-\min\{\frac1{3n},\frac3{16\kappaP}\}$ is
nondecreasing, $\rho_P\ge\rho_I$. The upper bound for $P=\hat\Sigma+\lambda I$ is
\eqref{eq:B1-kappaP-vs-kappa}.
\end{proof}

The mechanism is transparent: the supremum over $\theta$ in $R_P^2$ is attained on the decision boundary
of sample $i$ ($x_i^\top\theta=0$, curvature $\frac14x_ix_i^\top$), whereas the infimum over $\theta$ in
$\mu_P$ is approached far from all decision boundaries, where $F''\to\lambda I$ is isotropic; no single
metric can be adapted both to a rank-one curvature $\frac14x_ix_i^\top$ and to an isotropic one. (For the
square loss the curvature is constant, $F''\equiv\hat\Sigma+\lambda I=P$, every $f_i''=x_ix_i^\top+\lambda I$
has relative smoothness in $[n\,\mathrm{lev}_i,n\,\mathrm{lev}_i+1]$, and the global $\kappaP$ lies in
$[d_\infty(\lambda),d_\infty(\lambda)+1]$: a genuine gain. This is the ridge-regression setting of
\cite{gonen2016solving}, whose Theorem~4 bounds the \emph{average} condition number of the preconditioned
problem by $\tr(P^{-1/2}(\hat\Sigma+\lambda I)P^{-1/2})/\lmin(P^{-1/2}(\hat\Sigma+\lambda I)P^{-1/2})$,
equal to $d$ for $P=\hat\Sigma+\lambda I$; with uniform sampling and the maximum in
\eqref{eq:B1-RP-muP}, $d$ is replaced by $d_\infty(\lambda)$.) Hence, for logistic regression, any benefit
of the metric must come from the \emph{trajectory-dependent} constants, which we now quantify; this is
exactly the point of view of the B2 note.

\begin{proposition}[effective strong convexity from segment curvatures]\label{prop:B1-traj-sc}
For the logistic regression of Lemma~\ref{lem:B1-logistic}, every $\theta\ne\etastar$, $\delta:=\theta-\etastar$, and
$\bar\sigma_i'(\theta):=\int_0^1\sigma'\big(y_ix_i^\top(\etastar+s\delta)\big)\,ds\in(0,\tfrac14]$,
\begin{equation}\label{eq:B1-pairing-exact}
\fip(\theta)-\fip(\etastar)=\bar\sigma_i'(\theta)\,(x_i^\top\delta)\,x_i+\lambda\delta,\qquad
\inner{\Fp(\theta)}{\theta-\etastar}=\frac1n\sum_{i=1}^n\bar\sigma_i'(\theta)\,(x_i^\top\delta)^2+\lambda\norm{\delta}^2 .
\end{equation}
Consequently, with
$m(\theta):=\min_i\bar\sigma_i'(\theta)\ \ge\ \min_i\sigma'\big(\max\{|x_i^\top\etastar|,|x_i^\top\theta|\}\big)>0$:
\begin{enumerate}[label=(\roman*)]
\item for $P=\hat\Sigma+\lambda I$, the effective strong convexity
$\muEff^P(\theta):=\inner{\Fp(\theta)}{\delta}/\norm{\delta}_P^2$ is the convex combination
\begin{equation}\label{eq:B1-muEff-stat}
\muEff^P(\theta)=\frac{\frac1n\sum_i\bar\sigma_i'(\theta)(x_i^\top\delta)^2+\lambda\norm{\delta}^2}
{\frac1n\sum_i(x_i^\top\delta)^2+\lambda\norm{\delta}^2}\ \in\ [\,m(\theta),\,1\,]
\end{equation}
of the numbers $\bar\sigma_1'(\theta),\dots,\bar\sigma_n'(\theta),1$, and more precisely
\begin{equation}\label{ineq:P-trajsc}
\muEff^P(\theta)\ \ge\ \frac{m(\theta)\,\sigma_1+\lambda}{\sigma_1+\lambda}
\ \ge\ \max\Big\{m(\theta),\ \frac{\lambda}{\sigma_1+\lambda}\Big\}=\max\{m(\theta),\mu_P\};
\end{equation}
\item for $P=I$,
$\muEff(\theta)=\frac1n\sum_i\bar\sigma_i'(\theta)\,(x_i^\top\delta/\norm{\delta})^2+\lambda
\in[\,m(\theta)\sigma_d+\lambda,\ \sigma_1/4+\lambda\,]$.
\end{enumerate}
\end{proposition}
\begin{proof}
By the fundamental theorem of calculus applied to $s\mapsto\ell'(y_ix_i^\top(\etastar+s\delta))$,
$\ell'(y_ix_i^\top\theta)-\ell'(y_ix_i^\top\etastar)=\int_0^1\ell''(y_ix_i^\top(\etastar+s\delta))\,ds\cdot y_ix_i^\top\delta
=\bar\sigma_i'(\theta)\,y_ix_i^\top\delta$; multiplying by $y_ix_i$ ($y_i^2=1$) and adding $\lambda\delta$
gives the first identity in \eqref{eq:B1-pairing-exact} by \eqref{eq:B1-logistic-hessians}. Pairing with
$\delta$, averaging over $i$ and using $\Fp(\etastar)=0$ gives the second. Lower bound on $m(\theta)$:
$s\mapsto x_i^\top(\etastar+s\delta)$ is affine, so $|x_i^\top(\etastar+s\delta)|\le\max\{|x_i^\top\etastar|,|x_i^\top\theta|\}$
on $[0,1]$, and $\sigma'$ is decreasing in $|u|$ (Lemma~\ref{lem:B1-logistic}).
(i) $\norm{\delta}_P^2=\delta^\top\hat\Sigma\delta+\lambda\norm{\delta}^2=\frac1n\sum_i(x_i^\top\delta)^2+\lambda\norm{\delta}^2$,
which gives \eqref{eq:B1-muEff-stat}; it is a convex combination with (normalized) weights
$(x_i^\top\delta)^2/n$ and $\lambda\norm{\delta}^2$, whence $m(\theta)\le\muEff^P\le1$ because
$\bar\sigma_i'\le\frac14<1$. For \eqref{ineq:P-trajsc}: the numerator is $\ge\delta^\top(m\hat\Sigma+\lambda I)\delta$
with $m=m(\theta)$, and since $\hat\Sigma$ and $P$ are simultaneously diagonalizable,
$\delta^\top(m\hat\Sigma+\lambda I)\delta\ge\min_k\frac{m\sigma_k+\lambda}{\sigma_k+\lambda}\,\delta^\top P\delta
=\frac{m\sigma_1+\lambda}{\sigma_1+\lambda}\norm{\delta}_P^2$, because $s\mapsto\frac{ms+\lambda}{s+\lambda}$ is
nonincreasing on $[0,\infty)$ when $m\le1$ (its derivative is $\lambda(m-1)/(s+\lambda)^2$). Finally
$\frac{m\sigma_1+\lambda}{\sigma_1+\lambda}$ is a convex combination of $m$ and $1$ with weights
$\frac{\sigma_1}{\sigma_1+\lambda}$ and $\frac{\lambda}{\sigma_1+\lambda}$, hence it is $\ge m$ and
$\ge\frac{\lambda}{\sigma_1+\lambda}$.
(ii) Divide \eqref{eq:B1-pairing-exact} by $\norm{\delta}^2$; the bounds use $m\le\bar\sigma_i'\le\frac14$
and $\sigma_d\norm{\delta}^2\le\delta^\top\hat\Sigma\delta\le\sigma_1\norm{\delta}^2$.
\end{proof}

Note that \eqref{ineq:P-trajsc} is the segment-averaged, \emph{per-sample} version of the pointwise
inequality $F''(\theta')\succeq\min_i\sigma_i'(\theta')\,\hat\Sigma+\lambda I$: integrating the latter along
the segment $[\etastar,\theta]$ would only produce the constant $\int_0^1\min_i\sigma_i'(\etastar+s\delta)\,ds$,
which is $\le\min_j\int_0^1\sigma_j'(\etastar+s\delta)\,ds=m(\theta)$ (for each $j$,
$\min_i\sigma_i'\le\sigma_j'$ pointwise), so \eqref{ineq:P-trajsc}, which uses the minimum of the per-sample
segment averages, is at least as strong as the integrated pointwise bound. Its global version only gives
$\mu_P=\lambda/(\sigma_1+\lambda)$ because $\inf_\theta\min_i\sigma_i'(\theta)=0$. The experiments can test
\eqref{ineq:P-trajsc} directly (its tightness ratio is
$\frac{m(\theta_{t-1})\sigma_1+\lambda}{\sigma_1+\lambda}\big/\muEff^P(\theta_{t-1})$), and
\eqref{eq:B1-muEff-stat} shows that $\muEff^P(\theta_{t-1})$ is a direction-weighted \emph{average} of the per-sample
curvatures along the segment, typically much larger than the minimum $m$: a single sample with a huge
margin makes $m(\theta)$ tiny but has a negligible weight in \eqref{eq:B1-muEff-stat} unless the error
$\delta$ is aligned with it.

\begin{corollary}[leverage-score bound along trajectories]
\label{cor:B1-kappa-traj}
Let $\kappa_{\mathrm{eff}}^P(\theta):=R_P^2/\muEff^P(\theta)$ be the effective relative condition number
at $\theta$ (the quantity governing the effective contraction \eqref{eq:B1-effective-rate}). Then, for
$P=\hat\Sigma+\lambda I$ and $m(\theta)$ as in Proposition~\ref{prop:B1-traj-sc},
\begin{equation}\label{eq:B1-kappa-traj-stat}
\kappa^P_{\mathrm{eff}}(\theta)\ \le\ \Big(\frac{d_\infty(\lambda)}{4}+\frac{\lambda}{\sigma_d+\lambda}\Big)
\frac{\sigma_1+\lambda}{m(\theta)\sigma_1+\lambda}\ \le\ \frac{d_\infty(\lambda)/4+1}{m(\theta)}
\ \le\ \frac{C\,d_{\mathrm{eff}}(\lambda)/4+1}{m(\theta)},
\end{equation}
the last inequality under Assumption~\ref{ass:B1-coherence}, whereas for $P=I$ the same argument only gives
\begin{equation}\label{eq:B1-kappa-traj-I}
\frac{\max_i\norm{x_i}^2/4+\lambda}{\sigma_1/4+\lambda}\ \le\ \kappa_{\mathrm{eff}}(\theta)\ \le\
\frac{\max_i\norm{x_i}^2/4+\lambda}{m(\theta)\sigma_d+\lambda},
\qquad\text{with}\qquad \frac{\max_i\norm{x_i}^2/4+\lambda}{\sigma_1/4+\lambda}\ \ge\ \frac{\tr\hat\Sigma+4\lambda}{\sigma_1+4\lambda}.
\end{equation}
\end{corollary}
\begin{proof}
Divide \eqref{eq:B1-stat-RP} by \eqref{ineq:P-trajsc}, then use $\lambda/(\sigma_d+\lambda)\le1$ and
$\frac{\sigma_1+\lambda}{m\sigma_1+\lambda}\le\frac1m$ (as $m\le1$). For $P=I$, divide $R^2$ by the bounds
of Proposition~\ref{prop:B1-traj-sc}(ii), and use $\max_i\norm{x_i}^2\ge\tr\hat\Sigma$.
\end{proof}

In words: along the trajectory, the effective condition number in the statistical metric is at most a
leverage count, $d_\infty(\lambda)/4+1\le C\,d_{\mathrm{eff}}(\lambda)/4+1$, divided by the \emph{smallest} of the
per-sample curvatures of the logistic loss averaged on the segment $[\etastar,\theta]$ (the typical,
direction-weighted curvature of \eqref{eq:B1-muEff-stat} does not enter the bound); the eigenvalues
$\sigma_1,\dots,\sigma_d$ of $\hat\Sigma$ have disappeared, except through
$d_\infty(\lambda)\le C\,d_{\mathrm{eff}}(\lambda)$, with
$d_{\mathrm{eff}}(\lambda)=\sum_k\sigma_k/(\sigma_k+\lambda)\le d$, which is a \emph{count} of the directions
not killed by the regularization and is insensitive to the dynamic range of the spectrum. In the Euclidean
metric, by contrast, even the trajectory-dependent bound keeps the factor
$\max_i\norm{x_i}^2/(m\sigma_d+\lambda)\ge\tr\hat\Sigma/(\sigma_d/4+\lambda)$, the trace of $\hat\Sigma$
divided by its smallest eigenvalue (or by $\lambda$ when $\lambda\gg m\sigma_d$); and the lower bound in
\eqref{eq:B1-kappa-traj-I} shows that no trajectory can make the Euclidean $\kappa_{\mathrm{eff}}$ smaller
than $(\tr\hat\Sigma)/\sigma_1$ (up to the $\lambda$ terms), the stable rank $\norm{X}_F^2/\norm{X}_{\mathrm{op}}^2$ of $X$. The
price paid in the $P$-metric is the coherence constant $C$ and the curvature $m(\theta)$, which is small
when some margins $|x_i^\top\theta|$ are large ($\sigma'(u)\approx e^{-|u|}$); the Newton metric (d)
addresses this second issue.

On the instances of our experiments this price dominates, and Corollary~\ref{cor:B1-kappa-traj} predicts no
gain. In the experiment of Figure~B1 of the report ($n=1000$, $d=50$, $\kappa=10^3$, covariance spectrum
$j^{-\alpha}$; Section~\ref{sec:B1-predictions}), $m(\theta)\to\min_iw_i=7.4\cdot10^{-3}$ ($\alpha=1$) and
$6.1\cdot10^{-3}$ ($\alpha=2$) as $\theta\to\etastar$, with $w_i=\sigma_i'(\etastar)$ as in paragraph~(d), so
that the first bound in \eqref{eq:B1-kappa-traj-stat} tends to $1.7\cdot10^3$ and $1.2\cdot10^3$: above
$\kappa=10^3$, and above the upper bound in \eqref{eq:B1-kappa-traj-I} for $P=I$ with the same $m$ ($968$ and
$999$). Along the run of seed~$0$ it exceeds $\kappa$ from pass~$0.6$ on, and lies between $1.5\cdot10^3$ and
$1.8\cdot10^3$ ($\alpha=1$), resp.\ $1.1\cdot10^3$ and $1.3\cdot10^3$ ($\alpha=2$), over passes 2--20; the
condition \eqref{eq:B1-bottleneck-condition} below, even with $m=\min_iw_i$ (no lower bound on $m(\theta_t)$
along a trajectory converging to $\etastar$ can be larger), would require $n\ge4859$ and $n\ge3288$. By contrast,
$R_P^2/\lmin(P^{-1/2}F''(\etastar)P^{-1/2})$, which bounds $\limsup_{\theta\to\etastar}\kappa^P_{\mathrm{eff}}(\theta)$
by Lemma~\ref{lem:B1-asymptotic}, equals $166$ and $103$ (against $R^2/\lmin(F''(\etastar))=576$ and $981$ for
$P=I$), and the values of $\kappa^P_{\mathrm{eff}}(\theta_{t-1})$ measured over passes 15--20 are about $84$
and $31$ (against $460$ and $878$ for $P=I$; geometric means over states and seeds), $R_P^2$ being here the
upper bound of
\eqref{eq:B1-logistic-RP}, as in the experiments. The reason is that
$\lmin(P^{-1/2}F''(\etastar)P^{-1/2})=0.117$ and $0.109$ is the smallest over directions $v$ of a weighted
\emph{average} of $w_1,\dots,w_n,1$ (weights $(x_i^\top v)^2/n$ and $\lambda\norm{v}^2$), $16$ and $18$ times
$\min_iw_i$. The observed gain of $\hat\Sigma+\lambda I$ is therefore not certified by
Corollary~\ref{cor:B1-kappa-traj}; it is described only by the local constants at the optimum, which
Lemma~\ref{lem:B1-asymptotic} justifies as $\theta\to\etastar$ and which Section~\ref{sec:B1-predictions} uses
heuristically along the whole run, with the weighted average $\bar w$ in place of the minimum $m(\theta)$.

\begin{remark}[memory-bottleneck regime]\label{rem:B1-memory-bottleneck}
The one-step contraction factor is at most $1-\min\{\frac1{3n},\frac{3}{16\kappa'}\}$, with $\kappa'=\kappaP$
for the certified bound \eqref{eq:B1-onestep} and $\kappa'=R_P^2/\muEff^P(\theta_{t-1})$ for the state-dependent
bound \eqref{eq:B1-effective-rate} (a rate only under a deterministic bound on $\kappa'$ along the trajectory);
the memory term $\frac1{3n}$ is the binding one iff $\kappa'\le\frac{9n}{16}$. For
$P=\hat\Sigma+\lambda I$ and the trajectory constants of Corollary~\ref{cor:B1-kappa-traj}, a sufficient
condition is therefore
\begin{equation}\label{eq:B1-bottleneck-condition}
n\ \ge\ \frac{16}{9}\cdot\frac{d_\infty(\lambda)/4+1}{m}=\frac{4d_\infty(\lambda)+16}{9m}
\qquad\Big(\le\frac{4C\,d_{\mathrm{eff}}(\lambda)+16}{9m}\ \text{ under Assumption~\ref{ass:B1-coherence}}\Big),
\end{equation}
with $m$ a deterministic constant such that $m(\theta_t)\ge m$ for all $t$ almost surely: a
trajectory-dependent $m$ cannot be pulled out of the expectations when the one-step bounds are iterated (see
the case of a non-constant metric in Section~\ref{sec:B1-cost}), and without an a priori confinement of the
iterates Proposition~\ref{prop:B3-region} of the B3 note, transferred, gives the bound only up to the exit time
from a region (end of Section~\ref{sec:B1-choice}). In this regime ($n\gg d_{\mathrm{eff}}/m$) the bound on
the contraction of $P$-SAGA is $1-\frac1{3n}$ whatever $\lambda$ and the spectrum of $\hat\Sigma$: it needs
$3\log(2n/\varepsilon)$ effective passes (the $\log(2n)$ coming from the prefactor $1+\frac n4$ of
\eqref{eq:B1-rate}), which is the $n\log(1/\varepsilon)$ term of the lower bound of
hypothesis (A1). That bound is a worst case over the problem class for span-based methods with i.i.d.\ sampling,
which cover $P$-SAGA only through Proposition~\ref{prop:B1-transfer} and for $P$ fixed independently of the
data; with full-gradient steps this term becomes of order $\frac{n}{1+\ln(n/\kappa)}\log\frac1\varepsilon$ when $\kappa\ll n$
(A1 note), and the accelerated complexity $O((n+\sqrt{n\kappa})\log(1/\varepsilon))$ of (A3) keeps it. For $P=I$ the
corresponding condition $n\ge\frac{16}{9}\frac{\max_i\norm{x_i}^2/4+\lambda}{m\sigma_d+\lambda}$ fails by a
factor of order $\max_i\norm{x_i}^2/(d_\infty(\lambda)\,\sigma_d)$ when $\lambda\ll m\sigma_d$; and with
global constants both conditions read $n\ge\frac{16}{9}\kappaP\ge\frac{16}{9}\kappa$ by
Proposition~\ref{prop:B1-nofreelunch}. This fails, by at least the factor $\frac{16}9$ and for every $P$, when
$\lambda=R^2/n$, i.e.\ $\kappa=n$, which is the setting of the experiment of Section~\ref{sec:B1-predictions}.
The experiments of \cite[Sec.~5.4.4]{bach2024learning} take $\mu=R_x^2/n$ with $R_x$ the bound on
$\norm{x_i}$: this is $\kappa=n$ in the book's normalization (smoothness constant $R_x^2$), but
$\kappa=\frac n4+1\le\frac{9n}{16}$ (for $n\ge4$) with the exact constant $R^2=R_x^2/4+\lambda$ used here, so
that there the condition holds for $P=I$ and the memory term is already the binding one with global constants
(in our reproduction of that setting, Figure~0 of the report with $d=40$, it also holds for
$P=\hat\Sigma+\lambda I$: $\kappaP\le347$ for $n=10^3$ and $\kappaP\le2814$ for $n=10^4$).
\end{remark}

\begin{remark}[the variant $\frac14\hat\Sigma+\lambda I$]\label{rem:B1-quarter}
$P'=\frac14\hat\Sigma+\lambda I=\frac14(\hat\Sigma+4\lambda I)$ is the exact global upper envelope of
$F''$ ($F''(\theta)\preceq P'$ for all $\theta$ by \eqref{eq:B1-logistic-hessians}), so that $F$ is
$1$-smooth relative to $P'$. It is the statistical preconditioner $P_{4\lambda}:=\hat\Sigma+4\lambda I$ at
regularization $4\lambda$, rescaled by $\frac14$ (Remark~\ref{rem:B1-scaling}; the loss keeps its
regularization $\lambda$), so that $\kappa_{P'}=\kappa_{P_{4\lambda}}$ and
$\kappa^{P'}_{\mathrm{eff}}(\theta)=\kappa^{P_{4\lambda}}_{\mathrm{eff}}(\theta)$ (both $R_P^2$ and
$\muEff^P(\theta)$ are multiplied by $1/c$ under $P\mapsto cP$). By Lemma~\ref{lem:B1-leverage}(i) at
regularization $4\lambda$, $\norm{x_i}_{P_{4\lambda}^{-1}}^2=n\,\mathrm{lev}_i(4\lambda)$, so Lemma~\ref{lem:B1-logistic}
gives $R_{P_{4\lambda}}^2\le\frac{d_\infty(4\lambda)}4+\frac{\lambda}{\sigma_d+4\lambda}$ and
$\mu_{P_{4\lambda}}=\frac{\lambda}{\sigma_1+4\lambda}$; moreover $\muEff^{P'}(\theta)$ is a convex combination
of $4\bar\sigma_1',\dots,4\bar\sigma_n',1$ (weights $(x_i^\top\delta)^2/(4n)$ and $\lambda\norm{\delta}^2$),
hence lies in $[4m(\theta),1]$, and the argument of \eqref{ineq:P-trajsc}, now with the function
$s\mapsto\frac{ms+\lambda}{s+4\lambda}$ (nonincreasing since $m\le\frac14$), gives
$\muEff^{P_{4\lambda}}(\theta)\ge\frac{m(\theta)\sigma_1+\lambda}{\sigma_1+4\lambda}$. Hence
\begin{align}
\kappa_{P'}&\le\Big(\frac{d_\infty(4\lambda)}4+\frac{\lambda}{\sigma_d+4\lambda}\Big)\frac{\sigma_1+4\lambda}{\lambda},\nonumber\\
\kappa^{P'}_{\mathrm{eff}}(\theta)&\le\Big(\frac{d_\infty(4\lambda)}4+\frac{\lambda}{\sigma_d+4\lambda}\Big)
\frac{\sigma_1+4\lambda}{m(\theta)\sigma_1+\lambda}\le\frac{d_\infty(4\lambda)/4+1}{m(\theta)} ,\label{eq:B1-quarter-bounds}
\end{align}
the last step because $m(\sigma_1+4\lambda)\le m\sigma_1+\lambda$ when $m\le\frac14$. Since
$\lambda\mapsto\mathrm{lev}_i(\lambda)$ is nonincreasing (Lemma~\ref{lem:B1-leverage}(iii)),
$d_\infty(4\lambda)\le d_\infty(\lambda)$, and $\frac{\lambda}{\sigma_d+4\lambda}\le\frac{\lambda}{\sigma_d+\lambda}$:
the brackets in \eqref{eq:B1-quarter-bounds} are never larger than those of \eqref{eq:B1-stat-muP} and
\eqref{eq:B1-kappa-traj-stat}, but the factor $\sigma_1+4\lambda$ replaces $\sigma_1+\lambda$. Thus the
variant is never worse than $P=\hat\Sigma+\lambda I$ in the global and in the finer trajectory bounds
\emph{up to the factor} $(\sigma_1+4\lambda)/(\sigma_1+\lambda)\le4$, which is negligible when
$\lambda\ll\sigma_1$, and never worse at all in the crude bound $(d_\infty/4+1)/m$; the experiments may use
either.
\end{remark}

\paragraph{(d) $P=H:=F''(\etastar)$ (oracle Newton metric).}
Let $w_i:=\sigma_i'(\etastar)=\sigma'(y_ix_i^\top\etastar)\in(0,\frac14]$, $W:=\diag(w_1,\dots,w_n)$ and
$\hat\Sigma_w:=\frac1nX^\top WX$ (eigenvalues $\sigma_1^w\ge\dots\ge\sigma_d^w\ge0$), so that
$H=\hat\Sigma_w+\lambda I$ by \eqref{eq:B1-logistic-hessians}. Define the $w$-weighted ridge leverage
scores and effective dimension
\begin{align}
\mathrm{lev}^w_i(\lambda)&:=w_i\,x_i^\top(X^\top WX+n\lambda I)^{-1}x_i
=\big[W^{1/2}X(X^\top WX+n\lambda I)^{-1}X^\top W^{1/2}\big]_{ii},\label{eq:B1-wleverage}\\
d^w_{\mathrm{eff}}(\lambda)&:=\tr\big(\hat\Sigma_w(\hat\Sigma_w+\lambda I)^{-1}\big),\qquad
d^w_\infty(\lambda):=n\max_i\mathrm{lev}^w_i(\lambda),\nonumber
\end{align}
i.e. the ridge leverage scores of the weighted design $W^{1/2}X$ (rows $\sqrt{w_i}\,x_i^\top$).
Lemma~\ref{lem:B1-leverage} applied to $W^{1/2}X$ gives $\norm{\sqrt{w_i}\,x_i}_{H^{-1}}^2=n\,\mathrm{lev}^w_i(\lambda)$,
$\sum_i\mathrm{lev}^w_i=d^w_{\mathrm{eff}}$ and $0\le\mathrm{lev}^w_i<1$.

\begin{proposition}[constants in the Newton metric]\label{prop:B1-newton}
For $P=H=F''(\etastar)$:
\begin{enumerate}[label=(\roman*)]
\item $\norm{x_i}_{H^{-1}}^2=n\,\mathrm{lev}^w_i(\lambda)/w_i$, hence
\begin{equation}\label{eq:B1-newton-RP}
\max_i\frac{n\,\mathrm{lev}^w_i(\lambda)}{4w_i}\ \le\ R_H^2\ \le\ \max_i\frac{n\,\mathrm{lev}^w_i(\lambda)}{4w_i}+\frac{\lambda}{\sigma^w_d+\lambda}
\ \le\ \frac{d^w_\infty(\lambda)}{4\min_iw_i}+1,\qquad \mu_H=\frac{\lambda}{\sigma^w_1+\lambda}.
\end{equation}
\item For $\theta\neq\etastar$, $\delta:=\theta-\etastar$ and $\bar\sigma_i'(\theta)$ as in
Proposition~\ref{prop:B1-traj-sc},
\begin{equation}\label{eq:B1-muEff-newton}
\muEff^H(\theta)=\frac{\frac1n\sum_i\bar\sigma_i'(\theta)(x_i^\top\delta)^2+\lambda\norm{\delta}^2}
{\frac1n\sum_iw_i(x_i^\top\delta)^2+\lambda\norm{\delta}^2}
\ \in\ \Big[\min\Big\{1,\min_i\frac{\bar\sigma_i'(\theta)}{w_i}\Big\},\ \max\Big\{1,\max_i\frac{\bar\sigma_i'(\theta)}{w_i}\Big\}\Big],
\end{equation}
a convex combination of $\bar\sigma_1'/w_1,\dots,\bar\sigma_n'/w_n,1$; in particular
$\muEff^H(\theta)\to1$ as $\theta\to\etastar$, and $\muEff^H(\theta)\ge\mu_H$ always.
\end{enumerate}
\end{proposition}
\begin{proof}
(i) $x_i^\top H^{-1}x_i=w_i^{-1}(\sqrt{w_i}x_i)^\top H^{-1}(\sqrt{w_i}x_i)=n\,\mathrm{lev}^w_i/w_i$ by
Lemma~\ref{lem:B1-leverage}(i) for the design $W^{1/2}X$ (whose empirical covariance is $\hat\Sigma_w$);
then \eqref{eq:B1-logistic-RP}--\eqref{eq:B1-logistic-muP} with $\lmin(H)=\sigma^w_d+\lambda$,
$\lmax(H)=\sigma^w_1+\lambda$, and $\lambda/(\sigma_d^w+\lambda)\le1$.
(ii) Divide \eqref{eq:B1-pairing-exact} by $\norm{\delta}_H^2=\frac1n\sum_iw_i(x_i^\top\delta)^2+\lambda\norm{\delta}^2$;
the (normalized) weights are $w_i(x_i^\top\delta)^2/n$ for the values $\bar\sigma_i'/w_i$ and
$\lambda\norm{\delta}^2$ for the value $1$. As $\theta\to\etastar$, $\bar\sigma_i'(\theta)\to\sigma_i'(\etastar)=w_i$
by continuity of $\sigma'$, so every value tends to $1$. The last claim is \eqref{ineq:P-sc}.
\end{proof}

Thus the Newton metric makes the strong convexity side perfect near the optimum ($\muEff^H\to1$), but the
smoothness side still pays the supremum over $\theta$ of the per-sample curvature: sample $i$ enters
$R_H^2$ with the weight $\sup_\theta\sigma_i'(\theta)=\frac14$ in the numerator and with its curvature at
the optimum $w_i$ in the metric, hence the factor $1/w_i$ in \eqref{eq:B1-newton-RP}; a sample with a large
margin at $\etastar$ (small $w_i$) but a non-negligible weighted leverage is the worst case. Under a
\emph{homogeneous curvature} assumption $\min_iw_i\ge w_{\min}$ and coherence of the weighted leverage
scores ($d^w_\infty\le C_w\,d^w_{\mathrm{eff}}$ for a constant $C_w\ge1$), one gets
$R_H^2\le C_w\,d^w_{\mathrm{eff}}(\lambda)/(4w_{\min})+1$
and $\kappa_{\mathrm{eff}}^H(\theta)=R_H^2/\muEff^H(\theta)\to R_H^2\le C_w\,d^w_{\mathrm{eff}}/(4w_{\min})+1$
as $\theta\to\etastar$. The next lemma shows that this $1/w_{\min}$ is again an artefact of the global
supremum: the \emph{realized} smoothness constant $\LEff^H$ is leverage-controlled, with no curvature factor.

\begin{lemma}[limits of the effective constants at $\etastar$]\label{lem:B1-asymptotic}
Let the $f_i$ be $C^2$, $H_i:=f_i''(\etastar)$, $H:=F''(\etastar)=\frac1n\sum_iH_i\succ0$, and let
$\theta\to\etastar$ with $\delta:=\theta-\etastar$ satisfying $\delta/\norm{\delta}\to v$, $\norm{v}=1$. Then, for
every SPD $P$,
\begin{align}
\LEff^P(\theta)&\ \longrightarrow\ \frac{\frac1n\sum_i v^\top H_iP^{-1}H_iv}{v^\top Hv},\label{eq:B1-asymptotic}\\
\muEff^P(\theta)&\ \longrightarrow\ \frac{v^\top Hv}{v^\top Pv}
\in\big[\lmin(P^{-1/2}HP^{-1/2}),\lmax(P^{-1/2}HP^{-1/2})\big].\nonumber
\end{align}
Moreover $\frac1n\sum_iH_iP^{-1}H_i\succeq HP^{-1}H$, so the limit of $\LEff^P$ is at least
$\lmin(P^{-1/2}HP^{-1/2})$, and for $P=H$
\begin{equation}\label{eq:B1-asymptotic-newton}
\lim\LEff^H(\theta)=1+\frac{\frac1n\sum_i\norm{H^{-1/2}(H_i-H)v}^2}{\norm{H^{1/2}v}^2}\ \ge\ 1,
\end{equation}
with equality iff $H_iv=Hv$ for all $i$ (homogeneous per-sample curvature in the direction $v$).
For the logistic regression of Lemma~\ref{lem:B1-logistic}, in addition,
$\lim\LEff^H(\theta)\le\max\{d^w_\infty(\lambda),2\}$ for $P=H$,
$\lim\LEff^P(\theta)\le d_\infty(\lambda)/4+\frac32$ for $P=\hat\Sigma+\lambda I$, and, for $P=I$, with
$c_i:=x_i^\top v$ and $B_v:=\frac1n\sum_iw_i^2c_i^2\norm{x_i}^2$,
\begin{equation}\label{eq:B1-asymptotic-I}
\lim\LEff(\theta)=\frac{\lambda^2+2\lambda\,v^\top\hat\Sigma_wv+B_v}{\lambda+v^\top\hat\Sigma_wv},
\qquad
(v^\top\hat\Sigma_wv)^2\le B_v\le\frac{\max_i\norm{x_i}^2}{4}\,v^\top\hat\Sigma_wv=(\kappa-1)\lambda\,v^\top\hat\Sigma_wv
\end{equation}
(this is the exact formula $\LEff=(\lambda^2+2\lambda A+B)/(\lambda+A)$ of Lemma~\ref{lem:B2-logistic-formulas}, in
the limit, where its $A$ and $B$ become $v^\top\hat\Sigma_wv$ and $B_v$); in
particular $\lambda+v^\top\hat\Sigma_wv\le\lim\LEff(\theta)\le\lambda+(\kappa+1)\,v^\top\hat\Sigma_wv$.
\end{lemma}
\begin{proof}
Since $f_i$ is $C^2$, $\fip(\theta)-\fip(\etastar)=H_i\delta+o(\norm{\delta})$ and
$\Fp(\theta)=H\delta+o(\norm{\delta})$ (as $\Fp(\etastar)=0$). Hence
$\frac1n\sum_i\norm{\fip(\theta)-\fip(\etastar)}_{P^{-1}}^2=\frac1n\sum_i\delta^\top H_iP^{-1}H_i\delta+o(\norm{\delta}^2)$,
$\inner{\Fp(\theta)}{\delta}=\delta^\top H\delta+o(\norm{\delta}^2)$ and $\norm{\delta}_P^2=\delta^\top P\delta$;
dividing numerators and denominators by $\norm{\delta}^2$ and letting $\delta/\norm{\delta}\to v$ gives
\eqref{eq:B1-asymptotic}, because $v^\top Hv>0$ and $v^\top Pv>0$; the bracket is the range of the
Rayleigh quotient of $P^{-1/2}HP^{-1/2}$. The Loewner inequality:
$0\preceq\frac1n\sum_i(H_i-H)P^{-1}(H_i-H)=\frac1n\sum_iH_iP^{-1}H_i-HP^{-1}H$, since the cross terms
give $-\frac1n\sum_i(H_iP^{-1}H+HP^{-1}H_i)=-2HP^{-1}H$. Then
$v^\top HP^{-1}Hv=\norm{M\,P^{1/2}v}^2$ with $M:=P^{-1/2}HP^{-1/2}$, and for an SPD matrix $M$ and any $w$,
$\norm{Mw}^2=w^\top M^2w\ge\lmin(M)\,w^\top Mw$ (because $M^2-\lmin(M)M=M^{1/2}(M-\lmin(M)I)M^{1/2}\succeq0$),
with $w^\top Mw=v^\top Hv$ here; with $P=H$ the same decomposition gives \eqref{eq:B1-asymptotic-newton}.
For logistic regression, $H_i=w_ix_ix_i^\top+\lambda I$ and
$v^\top H_iP^{-1}H_iv=w_i^2(x_i^\top v)^2\norm{x_i}_{P^{-1}}^2+2\lambda w_i(x_i^\top v)(x_i^\top P^{-1}v)+\lambda^2v^\top P^{-1}v$.
For $P=H$: $w_i\norm{x_i}_{H^{-1}}^2=n\,\mathrm{lev}^w_i$, so the first terms average to
$\frac1n\sum_i n\,\mathrm{lev}^w_i\,w_i(x_i^\top v)^2\le d^w_\infty\,v^\top\hat\Sigma_wv$; the remaining terms average to
$2\lambda v^\top\hat\Sigma_wH^{-1}v+\lambda^2v^\top H^{-1}v=2\lambda v^\top(H-\lambda I)H^{-1}v+\lambda^2v^\top H^{-1}v
=2\lambda\norm{v}^2-\lambda^2v^\top H^{-1}v\le2\lambda\norm{v}^2$; hence the numerator is
$\le d^w_\infty\,v^\top\hat\Sigma_wv+2\lambda\norm{v}^2\le\max\{d^w_\infty,2\}\,v^\top Hv$.
For $P=\hat\Sigma+\lambda I$: $\norm{x_i}_{P^{-1}}^2=n\,\mathrm{lev}_i$ and $w_i\le\frac14$ give
$\frac1n\sum_iw_i^2(x_i^\top v)^2\,n\,\mathrm{lev}_i\le\frac{d_\infty}{4}\,v^\top\hat\Sigma_wv$; next
$P^{-1}\preceq\lambda^{-1}I$ gives $\lambda^2v^\top P^{-1}v\le\lambda\norm{v}^2$; finally, by
Cauchy--Schwarz, $\hat\Sigma_w\preceq\frac14\hat\Sigma$ and $\hat\Sigma\preceq P$,
\begin{align*}
2\lambda\,v^\top\hat\Sigma_wP^{-1}v&\le2\lambda\norm{\hat\Sigma_w^{1/2}v}\,\norm{\hat\Sigma_w^{1/2}P^{-1}v}
\le\lambda\norm{\hat\Sigma_w^{1/2}v}\,\norm{\hat\Sigma^{1/2}P^{-1}v}
\le\lambda\norm{\hat\Sigma_w^{1/2}v}\,\norm{P^{-1/2}v}\\
&\le\sqrt\lambda\,\norm{\hat\Sigma_w^{1/2}v}\,\norm{v}\le\tfrac12\big(v^\top\hat\Sigma_wv+\lambda\norm{v}^2\big),
\end{align*}
where we used $\norm{\hat\Sigma^{1/2}P^{-1}v}^2=v^\top P^{-1}\hat\Sigma P^{-1}v\le v^\top P^{-1}v=\norm{P^{-1/2}v}^2\le\norm{v}^2/\lambda$
and $ab\le\frac12(a^2+b^2)$. Summing the three bounds, the numerator is
$\le(\frac{d_\infty}4+\frac12)v^\top\hat\Sigma_wv+\frac32\lambda\norm{v}^2\le(\frac{d_\infty}4+\frac32)\,v^\top Hv$.
For $P=I$: $v^\top H_i^2v=\norm{H_iv}^2=\norm{w_ic_ix_i+\lambda v}^2=w_i^2c_i^2\norm{x_i}^2+2\lambda w_ic_i^2+\lambda^2$
(using $x_i^\top v=c_i$ and $\norm v=1$); averaging over $i$ gives $B_v+2\lambda v^\top\hat\Sigma_wv+\lambda^2$,
and $v^\top Hv=v^\top\hat\Sigma_wv+\lambda$, which is \eqref{eq:B1-asymptotic-I}. For the bounds on $B_v$:
$\norm{x_i}^2\ge c_i^2$ (Cauchy--Schwarz, $\norm v=1$) and the convexity of the square give
$B_v\ge\frac1n\sum_iw_i^2c_i^4\ge\big(\frac1n\sum_iw_ic_i^2\big)^2=(v^\top\hat\Sigma_wv)^2$, while
$w_i\le\frac14$ and $\norm{x_i}^2\le\max_j\norm{x_j}^2=4(\kappa-1)\lambda$ (by \eqref{eq:B1-kappa-I}) give
$B_v\le\frac{\max_j\norm{x_j}^2}4\cdot\frac1n\sum_iw_ic_i^2$. The two-sided bound on the limit follows, with
$A:=v^\top\hat\Sigma_wv$, from $(\lambda^2+2\lambda A+B_v)/(\lambda+A)\ge(\lambda+A)^2/(\lambda+A)$ and
$(\lambda^2+2\lambda A+B_v)/(\lambda+A)\le(\lambda^2+2\lambda A+(\kappa-1)\lambda A)/\lambda$.
\end{proof}

\paragraph{Global versus trajectory-dependent constants: where the supremum over $\theta$ is needed.}
Proposition~5.9 uses (I2) for each $i$ at the pair $(\theta_{t-1},\etastar)$ for every $t$, and (I6) at
$(\theta_0,\etastar)$. The proof of co-coercivity in Proposition~5.4 of \cite{bach2024learning} uses
smoothness at an auxiliary point, hence the global $R_P^2$; for $C^2$ convex $f_i$, however, the segment
suffices: with $\delta=\theta_{t-1}-\etastar$ and $\bar A_i:=\int_0^1f_i''(\etastar+s\delta)\,ds\succeq0$ one has
$\fip(\theta_{t-1})-\fip(\etastar)=\bar A_i\delta$ and, since $M^2\preceq\lmax(M)M$ for every $M\succeq0$,
\[
\norm{\bar A_i\delta}_{P^{-1}}^2\le\lmax\big(P^{-1/2}\bar A_iP^{-1/2}\big)\,
\inner{\fip(\theta_{t-1})-\fip(\etastar)}{\delta},
\]
where for logistic regression $\bar A_i=\bar\sigma_i'(\theta_{t-1})x_ix_i^\top+\lambda I$ (as in
\eqref{eq:B1-pairing-exact}), so that
$\lmax(P^{-1/2}\bar A_iP^{-1/2})\le\bar\sigma_i'(\theta_{t-1})\norm{x_i}_{P^{-1}}^2+\lambda/\lmin(P)$. The value
$\sup_\theta\sigma_i'=\frac14$ thus enters only as the worst case over trajectories (the same holds for (I6)):
the supremum over $\theta$ in $R_P^2$ is needed only because the step $\gamma=1/(4R_P^2)$ is fixed before the
trajectory is known. In the experiment of Section~\ref{sec:B1-predictions} the largest of these per-sample
segment constants at $\etastar$ is still $0.75R^2$ for $P=I$ and $0.85R_P^2$ for $P=\hat\Sigma+\lambda I$ at
$\alpha=1$ ($0.62R^2$ and $0.94R_P^2$ at $\alpha=2$, with $R_P^2$ the supremum of \eqref{eq:B1-RP-muP}), some
samples lying near their decision boundary: the slack of (I2) lies mostly in the average over samples and
directions that defines $\LEff^P$, not in the supremum over $\theta$. Strong convexity (I4) is also needed
only at the pairs $(\theta_{t-1},\etastar)$, and the exact formula \eqref{eq:B1-pairing-exact} shows that it
is a direction-weighted average of curvatures on the segment $[\etastar,\theta_{t-1}]$, which depends on
the trajectory. The global constants $(R_P^2,\mu_P)$ are therefore pessimistic in two different ways:
$\mu_P$ by the gap between $\inf_\theta$ and the segment averages (Propositions~\ref{prop:B1-traj-sc} and
\ref{prop:B1-newton}), and $R_P^2$ by the gap between $\max_i\sup_\theta$ and the realized $\LEff^P$
(Lemma~\ref{lem:B1-asymptotic}). The effective constants of the B2 note, written in the $P$-metric
(Definition~\ref{def:B1-ratios}), measure both gaps: $\muEff^P(\theta_{t-1})$ and $\LEff^P(\theta_{t-1})$;
the statements above predict that preconditioning moves the gap from one inequality to the other
(Section~\ref{sec:B1-predictions}). A rigorous way to use trajectory-dependent constants is
Proposition~\ref{prop:B3-region} of the B3 note, transferred by Proposition~\ref{prop:B1-transfer}: on a set
$\Theta\ni\etastar$, with $\inf_\Theta\muEff^P$ and $\sup_\Theta\LEff^P$ in place of $\mu_P$ and of the
co-coercivity constant (the step staying $1/(4R_P^2)$), it bounds $\E[V_T^P\ind_{\{\tau>T\}}]$, $\tau$ being the
exit time from $\Theta$, and it gives a rate if the iterates stay in $\Theta$ almost surely, which is not known
a priori for SAGA (it can be enforced by a projection onto $\Theta$, a proximal variant whose analysis differs
from Proposition~5.9). For non-quadratic losses \cite[Sec.~4.2]{hendrikx2020statistically} constrain their
problem to a ball $B(0,D)$ through an indicator function in the composite term and prove a uniform
concentration of Hessians over it, the ball being needed because the iterates depend on the sample that
defines the preconditioner.

\subsection{Cost, and the case of a non-constant metric}
\label{sec:B1-cost}

\paragraph{Per iteration and per pass.}
One iteration of $P$-SAGA computes one gradient $f'_{i(t)}(\theta_{t-1})$ ($O(d)$, or
$O(\operatorname{nnz}(x_i))$ for sparse data), updates the memory vector $z_{i(t)}$ and the running sum
$\sum_iz_i$ ($O(d)$), and applies $P^{-1}$ to $g_t$: $O(d)$ if $P$ is diagonal, $O(d^2)$ if $P$ is dense
and a Cholesky factor $P=LL^\top$ (or $P^{-1}$ itself) has been precomputed (two triangular solves).
Hence the cost of one iteration is $c_I=O(d)$, $c_{P_{\mathrm J}}=O(d)$, $c_{\hat\Sigma+\lambda I}=c_H=O(d^2)$, and one
effective pass ($n$ iterations) costs $O(nd)$ for $P\in\{I,P_{\mathrm J}\}$ and $O(nd^2)$ for a dense $P$.
Precomputation: $\diag(\hat\Sigma)$ costs $O(nd)$; $\hat\Sigma$ costs $O(nd^2)$ (one matrix product) and its
Cholesky factorization $O(d^3)$; $F''(\etastar)$ needs $\etastar$ (an oracle; in practice a Hessian at a
reference point, same cost $O(nd^2)$). Memory: $O(nd)$ for the $z_i$ as in Proposition~5.9 (the usual
implementation for linear models stores one scalar $\ell'(y_ix_i^\top\theta_{\tau_i})$ per sample and keeps
the regularizer outside the memory, $O(n)$ memory; this is a composite variant of SAGA
\cite{defazio2014saga} whose analysis differs slightly from Proposition~5.9), plus $O(d^2)$ for the factor
of a dense $P$. A dense $P$ also destroys sparsity ($P^{-1}x_i$ is dense), whereas a diagonal $P$ preserves
it.

\paragraph{Total complexity.}
By Corollary~\ref{cor:B1-rate}, the time to reach accuracy $\varepsilon$ is
\begin{equation}\label{eq:B1-total-cost}
T_P\ \asymp\ c_P\Big(3n+\frac{16}{3}\kappaP\Big)\log\frac{2n}{\varepsilon}\ +\ \text{(precomputation)},
\end{equation}
to be compared with $T_I\asymp d\,(3n+\frac{16}{3}\kappa)\log\frac{2n}{\varepsilon}$ for SAGA
\cite[Sec.~5.4.4]{bach2024learning}. The diagonal metric is free. A dense metric multiplies the
per-iteration cost by $d$ and adds $O(nd^2+d^3)$, so it pays off only if the number of iterations
decreases by more than a factor $d$: dividing \eqref{eq:B1-total-cost} by $d\log\frac{2n}{\varepsilon}$, the
condition is
\[
d\Big(3n+\frac{16}{3}\kappaP\Big)+O\Big(\frac{nd+d^2}{\log(2n/\varepsilon)}\Big)\ \lesssim\ 3n+\frac{16}{3}\kappa,
\qquad\text{i.e. roughly}\qquad \kappa\ \gtrsim\ d\,(n+\kappaP),
\]
the very ill-conditioned regime $\kappa\gg nd$ (the regime $\kappa\gg n$ of hypothesis (A4), with an extra
factor $d$). With global constants this never happens for logistic regression ($\kappaP\ge\kappa$,
Proposition~\ref{prop:B1-nofreelunch}). With trajectory constants, \emph{both} condition numbers must be
replaced by their effective values: inserting the upper bounds \eqref{eq:B1-kappa-traj-stat} for
$P=\hat\Sigma+\lambda I$ and \eqref{eq:B1-kappa-traj-I} for $P=I$ (with the same lower bound $m$ on the
segment curvatures along both trajectories), the condition reads
\begin{equation}\label{eq:B1-breakeven}
\frac{\max_i\norm{x_i}^2/4+\lambda}{m\sigma_d+\lambda}\ \gtrsim\ d\,\Big(n+\frac{d_\infty(\lambda)/4+1}{m}\Big)
\end{equation}
if both bounds are tight. Since the left-hand side of \eqref{eq:B1-breakeven} is at most the global
$\kappa=\max_i\norm{x_i}^2/(4\lambda)+1$, replacing it by $\kappa$ (global-constant SAGA against
effective-constant $P$-SAGA, an asymmetric comparison) gives a condition,
$\kappa\gtrsim d(n+(d_\infty(\lambda)/4+1)/m)$, which is only \emph{necessary}: it is a lower bound on the
ill-conditioning required for a dense metric to pay off, and it overstates the benefit of the metric when
$\lambda\ll m\sigma_d$ (then the left-hand side of \eqref{eq:B1-breakeven} is of order
$\max_i\norm{x_i}^2/(4m\sigma_d)\ll\kappa$). In the memory-bottleneck regime of
Remark~\ref{rem:B1-memory-bottleneck}, $T_P\asymp3c_Pn\log(2n/\varepsilon)$: $3\log(2n/\varepsilon)$
passes, each $d$ times more expensive than a SAGA pass when $P$ is dense.

\paragraph{A non-constant metric breaks the transfer argument.}
Suppose the step uses $P_t$ at iteration $t$: $\theta_t=\theta_{t-1}-\gamma_tP_t^{-1}g_t$ (variable-metric
or quasi-Newton variance reduction). Proposition~\ref{prop:B1-transfer} no longer applies: the transformed
functions $\tilde f_i^{(t)}=f_i\circ P_t^{-1/2}$ and the transformed memory $P_t^{-1/2}z_i$ change with
$t$, so there is no single Euclidean SAGA run to which Proposition~5.9 could be applied. In the $P$-metric
proof of Section~\ref{sec:B1-proofP}, the one-step inequality \eqref{eq:B1-onestep} still holds at
iteration $t$ for the Lyapunov function $V_t^{P_t}$ built with $P_t$ and $\gamma_t=1/(4R_{P_t}^2)$,
\emph{provided} $P_t$ (hence $\gamma_t$, $R_{P_t}^2$, $\rho_{P_t}$) is measurable with respect to the past
$\mathcal F_{t-1}:=\sigma(i(1),\dots,i(t-1))$, i.e. is computed before $i(t)$ is drawn: conditionally on
$\mathcal F_{t-1}$ the metric is then fixed, the unbiasedness $\E[g_t\mid\mathcal F_{t-1}]=\Fp(\theta_{t-1})$
can be used inside the $P_t$-norms, and the proof applies verbatim, giving
$\E[V_t^{P_t}\mid\mathcal F_{t-1}]\le\rho_{P_t}V_{t-1}^{P_t}$ (if $P_t$ depended on $i(t)$, then
$\E[P_t^{-1}g_t\mid\mathcal F_{t-1}]\neq P_t^{-1}\Fp(\theta_{t-1})$ in general and the argument already
breaks at \eqref{eq:B1-expand}). But this compares $\E[V_t^{P_t}\mid\mathcal F_{t-1}]$ with
$V_{t-1}^{P_t}$, not with $V_{t-1}^{P_{t-1}}$; one must pay the change of metric
\begin{equation}\label{eq:B1-metric-drift}
V_{t-1}^{P_t}\le c_t\,V_{t-1}^{P_{t-1}},\qquad
c_t:=\max\Big\{\lmax\big(P_{t-1}^{-1/2}P_tP_{t-1}^{-1/2}\big),\
\frac{\gamma_t^2}{\gamma_{t-1}^2}\lmax\big(P_t^{-1/2}P_{t-1}P_t^{-1/2}\big)\Big\}
\end{equation}
(Lemma~\ref{lem:B1-geometry}(v) applied to the pairs $(P_{t-1},P_t)$ and $(P_{t-1}^{-1},P_t^{-1})$; with
the convention $P_0:=P_1$, $\gamma_0:=\gamma_1$, so that $c_1=1$). Both $\rho_{P_t}$ and $c_t$ are
$\mathcal F_{t-1}$-measurable but random for an adaptive metric, and correlated with $V_{t-1}^{P_{t-1}}$,
so they cannot be pulled out of the expectation; if however $\rho_{P_s}c_s\le\bar\rho_s$ almost surely for
deterministic constants $\bar\rho_s$, then taking expectations and iterating gives
$\E V_t^{P_t}\le\prod_{s=1}^t\bar\rho_s\,V_0^{P_0}$. Within this Lyapunov argument, linear convergence is
thus certified only if the metric drift is controlled by deterministic bounds, e.g. $c_t\le\bar c_t$ with
$\sum_t\log\bar c_t<\infty$
(a stationary limit metric, reached fast enough), or
$c_t\le1+\frac12\min\{\frac1{3n},\frac3{16\bar\kappa}\}$ for a deterministic $\bar\kappa\ge\sup_t\kappa_{P_t}$
(then $\rho_{P_t}c_t\le1-\frac12\min\{\frac1{3n},\frac3{16\bar\kappa}\}$), which is the analogue of the
condition ``$\eta_t$ does
not vary too much (not more than a factor $1-O(1/n)$)'' stated after Theorem~4 of \cite{dragomir2021fast}
(there the gains $G_t$ must be nonincreasing for the Bregman potential to be a supermartingale). A uniform
sandwich $\gamma I\preceq P_t^{-1}\preceq\Gamma I$ alone gives only $c_t\le\Gamma/\gamma$ for a constant step
(from here on, $\gamma,\Gamma,\eta,\lambda,\Lambda,m,\mu$ in the quoted results are the symbols of the cited papers),
useless per step;
this is why the stochastic quasi-Newton variance-reduced methods of \cite{moritz2016linearly} (stochastic
L-BFGS with SVRG gradients; Lemma~4: $\gamma I\preceq H_r\preceq\Gamma I$) and \cite{gower2016stochastic}
(stochastic block BFGS with SVRG gradients; Assumption~1 and Lemma~1: $\gamma I\preceq H_t\preceq\Gamma I$)
are analysed in the Euclidean norm, the metric entering only through $(\gamma,\Gamma)$: Theorem~7 of
\cite{moritz2016linearly} gives, for $\lambda I\preceq\nabla^2f_T\preceq\Lambda I$, step $\eta$ and
inner-loop length $m$, the rate $\alpha=\big(\frac1{2m\eta}+\eta\Gamma^2\Lambda^2\big)\big/\big(\gamma\lambda-\eta\Gamma^2\Lambda^2\big)$,
and Theorem~1 of \cite{gower2016stochastic} the rate
$\rho=\big(\frac1{2m\eta}+\eta\Gamma^2\Lambda(\Lambda-\lambda)\big)\big/\big(\gamma\lambda-\eta\Gamma^2\Lambda^2\big)$.
In the first formula, $\alpha<1$ iff $\frac1{2m\eta}+2\eta\Gamma^2\Lambda^2<\gamma\lambda$; minimizing the
left-hand side over $\eta$ (AM--GM: its minimum is $2\Gamma\Lambda/\sqrt m$, attained at
$\eta=1/(2\Gamma\Lambda\sqrt m)$) shows that $\alpha<1$ for some $\eta$ if and only if
$m>4\big(\Gamma\Lambda/(\gamma\lambda)\big)^2\ge4(\Lambda/\lambda)^2$, since $\Gamma\ge\gamma$: the
certified epoch length is at least $4(\Lambda/\lambda)^2$. Since both papers require
$\lambda I\preceq\nabla^2f_T\preceq\Lambda I$ for \emph{every} nonempty subset $T$ of indices (single
components and the full sum included, $f_T$ being the average of the $f_i$, $i\in T$), $\Lambda\ge R^2$ and
$\lambda\le\mu$ with our $R^2$ and $\mu$, so this epoch length is at least $4\kappa^2$, whereas the epoch length
of Euclidean SVRG is linear in the condition number; the benefit of the metric is not captured at all.
Capturing it requires a relative-smoothness analysis with a fixed (or
asymptotically fixed) reference function, as here (\cite{dragomir2021fast}: Corollary~1 for quadratic $h$,
Theorem~4 with gains $G_t$ otherwise), or a preconditioner fixed in advance from statistics of the data: a
sub-sampled Hessian plus ridge as in \cite{hendrikx2020statistically}, whose Theorem~2 shows that for
quadratic losses with $\norm{a_i}\le R$, a sub-sample of size $n>\frac{28}{3}\log\frac{2d}{\delta}$ and the
ridge $\mu=\frac12\big(\frac{28R^2}{3n}\log\frac{2d}\delta-\lambda\big)^+$ give, with probability $1-\delta$,
$(\frac32+\frac{2\mu}{\lambda})^{-1}(H_f+\mu I)\preceq H_F\preceq2(H_f+\mu I)$, i.e. a relative condition
number $O\big(1+\frac{\kappa_\ell}{n}\log\frac d\delta\big)$ with $\kappa_\ell=R^2/\lambda$; or a sketched
preconditioner as in \cite{gonen2016solving} (rank-$k$ sketch of the covariance for ridge regression,
computed in time $O(ndk\log n)$; their Theorem~1 replaces the average condition number
$\sum_i(\sigma_i+\lambda)/(\sigma_d+\lambda)$ of SVRG by one of order $(k\sigma_k+\sum_{i>k}\sigma_i)/\lambda$).
In all these cases the reference function is fixed in advance; the metric itself is constant only for a
quadratic $h$ (quadratic losses in \cite{hendrikx2020statistically,gonen2016solving}, Corollary~1 of
\cite{dragomir2021fast}), which is exactly what Proposition~\ref{prop:B1-transfer} needs. (Non-uniform sampling
of $i(t)$ is outside the uniform-sampling setting of Proposition~5.9 and of this note. For the record: Gower
et al.\ \cite[Thm.~5.2 and Table~1, row~8]{gower2021stochastic} proved (see also \cite{qian2019saga}) that
SAGA with the reweighted unbiased estimator
$\frac1n\sum_iz_i^{(t-1)}+\frac1{np_{i(t)}}\big(f'_{i(t)}(\theta_{t-1})-z_{i(t)}^{(t-1)}\big)$ and the
probabilities $p_i\propto\mu n+4L_i$ that minimize their bound reaches accuracy $\varepsilon$ in
$O\big((n+4\bar L/\mu)\log\frac1\varepsilon\big)$ iterations, $\bar L=\frac1n\sum_iL_i$ being the
\emph{average} of the per-sample smoothness constants; this resolves what they call a conjecture of Schmidt
et al.\ \cite{schmidt2015nonuniform}, who had obtained a rate depending on $\bar L$ only for a variant drawing
two indices per iteration (one with $p_i\propto L_i$ for the gradient, one uniform for the memory update)
\cite[Prop.~1]{schmidt2015nonuniform}; a non-rigorous motivation of Lipschitz sampling for SAG is given in
\cite[Sec.~4.8]{schmidt2017minimizing}. The change of variables of
Proposition~\ref{prop:B1-transfer} applies verbatim to this variant, and for $P=\hat\Sigma+\lambda I$ the
per-sample relative smoothness constants satisfy $L_i\le\frac{n\,\mathrm{lev}_i(\lambda)}4+\frac{\lambda}{\sigma_d+\lambda}$
by the Weyl bound in the proof of Lemma~\ref{lem:B1-logistic}, so that $d_\infty(\lambda)$ would be replaced
by $d_{\mathrm{eff}}(\lambda)=\sum_i\mathrm{lev}_i(\lambda)$ in the \emph{global} bound \eqref{eq:B1-stat-muP};
whether the same substitution holds in the trajectory bounds of Corollary~\ref{cor:B1-kappa-traj} is not
covered by those analyses and is left open.)

\subsection{Predictions for the experiments}
\label{sec:B1-predictions}

Tables~\ref{tab:B1-predictions-global} and \ref{tab:B1-predictions-late} collect, for the four metrics,
the certified constants (Lemma~\ref{lem:B1-logistic}, Propositions~\ref{prop:B1-diag},
\ref{prop:B1-stat-global} and \ref{prop:B1-newton}), the certified rate
$\rho_P=1-\min\{\frac1{3n},\frac{3}{16\kappaP}\}$, and the predicted late-phase behaviour of the
effective constants (Lemma~\ref{lem:B1-asymptotic}). The late-phase entries are evaluated for the extreme
direction $v$ = bottom eigenvector of the relative Hessian $M_P:=P^{-1/2}F''(\etastar)P^{-1/2}$ (the most
slowly contracted direction), as if the late-phase error $\delta$ aligned with it, as it does for gradient
descent on a quadratic; $\bar w$ denotes a leverage-weighted average of the curvatures $w_i=\sigma_i'(\etastar)$
in that direction and $\overline{n\,\mathrm{lev}}$ a weighted average of the $n\,\mathrm{lev}_i(\lambda)$ (equal to
$d_{\mathrm{eff}}(\lambda)$ when leverage scores are homogeneous). Some entries do not depend on this
alignment. By Lemma~\ref{lem:B1-asymptotic} and the compactness of the unit sphere, whatever the direction of
the error, $\liminf_{\theta\to\etastar}\muEff^P(\theta)\ge\lmin(M_P)$, hence $\limsup r_{\mathrm{sc}}\le\mu_P/\lmin(M_P)$
and $\limsup R_P^2/\muEff^P\le R_P^2/\lmin(M_P)$; the bounds $\lim\LEff^P\le d_\infty(\lambda)/4+\frac32$
($P=\hat\Sigma+\lambda I$) and $\lim\LEff^H\in[1,\max\{d^w_\infty(\lambda),2\}]$ hold along every limiting
direction; and for $P=I$ the formulas for $\LEff$ and $r_{\mathrm{coco}}$ hold with $\sigma^w_d$ replaced by
$v^\top\hat\Sigma_wv\in[\sigma^w_d,\sigma^w_1]$. The other entries, in particular the approximate values
($\approx$), assume the alignment, a heuristic borrowed from gradient descent. For SAGA, the (heuristic)
criterion of Remark~\ref{rem:B2-P2-outcome} of the B2 note, transferred by Proposition~\ref{prop:B1-transfer},
expects alignment only when the parameter term is locally the active one, $9n\lmin(M_P)/(16R_P^2)<1$. On the
instances of Figure~B1 of the report this quantity is $0.98$, $0.88$, $3.4$, $6.8$ ($\alpha=1$) and $0.57$,
$0.23$, $5.5$, $13.6$ ($\alpha=2$) for $I$, $P_{\mathrm J}$, $\hat\Sigma+\lambda I$, $H$ (with $R_P^2$ the upper
bound of \eqref{eq:B1-logistic-RP} used in the experiments): the dense metrics are in the memory-bottleneck
regime of (P3), where the error stays spread (over passes 15--20, $\muEff^P/\lmin(M_P)\approx2.0$ and $3.3$ for
$\hat\Sigma+\lambda I$). Even for $P=I$ alignment is not observed at $\kappa=10^3$: in the B2 runs ($P=I$, same
problems, one seed) the final error has $0.31$ ($\alpha=1$) and $0.0025$ ($\alpha=2$) of its energy on the bottom
eigenvector of $H$ (participation ratios $3.8$ and $12$), the bottom of the spectrum of $H$ being nearly
degenerate ($5$ and $29$ eigenvalues within $10\%$ of $\lmin(H)$); alignment occurs there only at $\kappa=10^4$
with $\alpha\in\{0,1\}$. For $P=I$ and $\alpha=2$, for instance, the limit of $r_{\mathrm{coco}}$ along the bottom
eigenvector is $0.0028$, against $0.013$ measured over passes 15--20.

\begin{table}[ht]
\centering
\footnotesize
\renewcommand{\arraystretch}{1.3}
\begin{tabular}{@{}p{2.0cm}p{2.6cm}p{3.2cm}p{3.4cm}p{3.4cm}@{}}
\toprule
 & $P=I$ & $P=P_{\mathrm J}=\diag(\hat\Sigma)+\lambda I$ & $P=\hat\Sigma+\lambda I$ & $P=H=F''(\etastar)$ \\
\midrule
$R_P^2$ (certified) & $\frac14\max_i\norm{x_i}^2+\lambda$
 & $\frac14\max_i\norm{x_i}_{P_{\mathrm J}^{-1}}^2+\frac{\lambda}{\min_j\hat\Sigma_{jj}+\lambda}$
 & $\frac{d_\infty(\lambda)}4+\frac{\lambda}{\sigma_d+\lambda}$ \newline $\le\frac{Cd_{\mathrm{eff}}(\lambda)}4+1$
 & $\max_i\frac{n\,\mathrm{lev}^w_i(\lambda)}{4w_i}+\frac{\lambda}{\sigma^w_d+\lambda}$ \\[2pt]
$\mu_P$ (global, exact) & $\lambda$ & $\frac{\lambda}{\max_j\hat\Sigma_{jj}+\lambda}$
 & $\frac{\lambda}{\sigma_1+\lambda}$ & $\frac{\lambda}{\sigma^w_1+\lambda}$ \\[2pt]
$\kappaP$ (global) & $\kappa$ & $\ge\kappa$ & $\in[\kappa,\ \kappa\operatorname{cond}(P)]$ & $\ge\kappa$ \\[2pt]
certified $\rho_P$ & $1-\min\{\frac1{3n},\frac3{16\kappa}\}$ & $\ge\rho_I$: not better than $P=I$ & $\ge\rho_I$: not better than $P=I$ & $\ge\rho_I$: not better than $P=I$ \\[2pt]
cost per pass & $O(nd)$ & $O(nd)$ & $O(nd^2)$ $+O(nd^2+d^3)$ once & $O(nd^2)$ $+O(nd^2+d^3)$ once, needs $\etastar$ \\
\bottomrule
\end{tabular}
\caption{Certified (global) constants for $\ell_2$-regularized logistic regression.}
\label{tab:B1-predictions-global}
\end{table}

\begin{table}[ht]
\centering
\footnotesize
\renewcommand{\arraystretch}{1.3}
\begin{tabular}{@{}p{1.6cm}p{3.2cm}p{2.8cm}p{3.6cm}p{3.4cm}@{}}
\toprule
late phase & $P=I$ & $P=P_{\mathrm J}$ & $P=\hat\Sigma+\lambda I$ & $P=H=F''(\etastar)$ \\
\midrule
$\muEff^P\to$ & $\lmin(H)=\sigma^w_d+\lambda$ & $\lmin(M_{P_{\mathrm J}})$
 & $\lmin(M_P)\approx\bar w$, \newline $\in\big[\frac{w_{\min}\sigma_1+\lambda}{\sigma_1+\lambda},1\big]$ & $1$ \\[2pt]
$\LEff^P\to$ & $\frac{\lambda^2+2\lambda\sigma^w_d+B_v}{\lambda+\sigma^w_d}$ \newline $\in[\sigma^w_d+\lambda,\ \lambda+(\kappa+1)\sigma^w_d]$ & same with $\norm{\cdot}_{P_{\mathrm J}^{-1}}$
 & $\le\frac{d_\infty(\lambda)}4+\frac32$, $\approx\bar w\,\overline{n\,\mathrm{lev}}$ & $\in[1,\max\{d^w_\infty(\lambda),2\}]$ \\[2pt]
$r_{\mathrm{sc}}\to$ & $\frac{\lambda}{\sigma^w_d+\lambda}$ ($\approx1$ if $\sigma^w_d\ll\lambda$) & $\mu_{P_{\mathrm J}}/\lmin(M_{P_{\mathrm J}})$
 & $\approx\frac{\lambda}{(\sigma_1+\lambda)\bar w}\ll1$ & $\frac{\lambda}{\sigma^w_1+\lambda}\ll1$ \\[5pt]
$r_{\mathrm{coco}}\to$ & $\frac{\lambda^2+2\lambda\sigma^w_d+B_v}{(\lambda+\sigma^w_d)R^2}$ \newline $\in\big[\frac1\kappa,\ \frac1\kappa+\frac{2\sigma^w_d}{\lambda}\big]$ \newline ($\approx\frac1\kappa$ if $\sigma^w_d\ll\lambda/\kappa$) & no prediction
 & $\approx4\bar w\,\overline{n\,\mathrm{lev}}/d_\infty(\lambda)\le1$ & $\frac{1+\text{heterogeneity}}{R_H^2}$ \\[2pt]
loosest step & (I2) co-coercivity & (I2) and (I4) & (I4) strong convexity & (I4) strong convexity \\[2pt]
$R_P^2/\muEff^P\to$ & $\frac{\max_i\norm{x_i}^2/4+\lambda}{\sigma^w_d+\lambda}$
 & {\raggedright within a factor $\operatorname{cond}(C_{\mathrm{corr}})$ of the $\hat\Sigma+\lambda I$ value
 (either direction)\par}
 & $\lesssim\frac{Cd_{\mathrm{eff}}(\lambda)/4+1}{\bar w}$ & $R_H^2\le\frac{C_wd^w_{\mathrm{eff}}(\lambda)}{4w_{\min}}+1$ \\
\bottomrule
\end{tabular}
\caption{Predicted late-phase behaviour ($M_P=P^{-1/2}F''(\etastar)P^{-1/2}$, $v$ the bottom eigenvector of
$M_P$, taken heuristically as the late-phase error direction (see the caveat before
Table~\ref{tab:B1-predictions-global}); the $\muEff^P$ row is also a lower bound for the $\liminf$, and the
$r_{\mathrm{sc}}$ and $R_P^2/\muEff^P$ rows upper bounds for the $\limsup$, whatever the direction of the error;
for $P=I$, $v^\top\hat\Sigma_wv=\sigma^w_d$ and
$B_v=\frac1n\sum_iw_i^2(x_i^\top v)^2\norm{x_i}^2\le(\kappa-1)\lambda\sigma^w_d$ as in
\eqref{eq:B1-asymptotic-I}; ``heterogeneity'' is the second term of \eqref{eq:B1-asymptotic-newton}).}
\label{tab:B1-predictions-late}
\end{table}

In words, with the statements they follow from:
\begin{enumerate}[label=(P\arabic*)]
\item \emph{Certified rates.} With global constants, none of the three preconditioners improves $\rho_P$
for logistic regression ($\rho_P\ge\rho_I$ since $\kappaP\ge\kappa$, Proposition~\ref{prop:B1-nofreelunch});
any observed speed-up must show up in the
effective constants, not in $\rho_P$. For the square loss the statistical preconditioner does improve
$\rho_P$, to at most $1-\min\{\frac1{3n},\frac{3}{16(d_\infty(\lambda)+1)}\}$.
\item \emph{Looseness moves from (I2) to (I4).} Under $P=I$ the late-phase error lives in the
low-curvature directions of $H$, where (I4) is nearly tight
($r_{\mathrm{sc}}\to\lambda/\lmin(H)=\lambda/(\sigma^w_d+\lambda)$ along the bottom eigenvector of $H$, and
$\limsup r_{\mathrm{sc}}\le\lambda/\lmin(H)$ in general; close to $1$ when $\sigma^w_d\ll\lambda$)
and (I2) is loose: by \eqref{eq:B1-asymptotic-I} with $v$ the bottom eigenvector of $H$ (so
$v^\top\hat\Sigma_wv=\sigma^w_d$),
\[
r_{\mathrm{coco}}\ \to\ \frac{\lambda^2+2\lambda\sigma^w_d+B_v}{(\lambda+\sigma^w_d)R^2}
\ \in\ \Big[\frac{\sigma^w_d+\lambda}{R^2},\ \frac1\kappa+\Big(1+\frac1\kappa\Big)\frac{\sigma^w_d}{\lambda}\Big],
\]
which is $\approx\lambda/R^2=1/\kappa$ when $\sigma^w_d\ll\lambda$ \emph{and} $B_v\ll\lambda^2$ (sufficient:
$\sigma^w_d\ll\lambda/\kappa$, since $B_v\le(\kappa-1)\lambda\sigma^w_d$), and in any case $\ll1$ as soon as
$\sigma^w_d\ll\lambda$. The condition $\sigma^w_d\ll\lambda$ alone does not give $\LEff\approx\lambda$: it
kills the term $2\lambda\sigma^w_d$ but not $B_v$, which may be as large as
$(\kappa-1)\lambda\sigma^w_d\gg\lambda^2$ on ill-conditioned logistic instances ($\kappa\gg n$); the experiments should
therefore record $B_v$ (or directly the limit value) rather than test $r_{\mathrm{coco}}\approx1/\kappa$.
Under $P=\hat\Sigma+\lambda I$ and $P=H$ the bottom eigenspace
of $M_P$ is a high-variance data direction with relative curvature $\approx\bar w$ (resp.\ $1$), so (I4)
with the global $\mu_P$ becomes the loosest step ($r_{\mathrm{sc}}\approx\lambda/((\sigma_1+\lambda)\bar w)$,
resp.\ $\lambda/(\sigma_1^w+\lambda)$), while $r_{\mathrm{coco}}$ rises to $\approx4\bar w\,\overline{n\,\mathrm{lev}}/d_\infty(\lambda)$
(Lemma~\ref{lem:B1-asymptotic}): it moves closer to $1$ when the per-sample curvature is homogeneous and
close to its maximum $\frac14$ (data near the decision boundary) \emph{and} the leverage scores are
homogeneous ($\overline{n\,\mathrm{lev}}\approx d_{\mathrm{eff}}\approx d_\infty$); it remains bounded by $4\bar w\le1$.
\item \emph{Effective rate.} By \eqref{eq:B1-effective-rate} the parameter part contracts at
$1-(1-\frac58r_{\mathrm{coco}})\muEff^P/(2R_P^2)$; the predicted late-phase effective condition numbers
$R_P^2/\muEff^P$ are the last row of Table~\ref{tab:B1-predictions-late}: equal to
$(\max_i\norm{x_i}^2/4+\lambda)/(\sigma^w_d+\lambda)=R^2/\lmin(F''(\etastar))\ \ge\ (\tr\hat\Sigma/4+\lambda)/(\sigma^w_d+\lambda)$
for $P=I$ (the anisotropy of $\hat\Sigma$), $\lesssim(Cd_{\mathrm{eff}}/4+1)/\bar w$ for the statistical metric
(no anisotropy), $\to R_H^2$ for the Newton metric; for the Jacobi metric Proposition~\ref{prop:B1-diag}(iii)
only gives
$\operatorname{cond}(C_{\mathrm{corr}})^{-1}\le\kappa^{P_{\mathrm J}}_{\mathrm{eff}}/\kappa^{\hat\Sigma+\lambda I}_{\mathrm{eff}}\le\operatorname{cond}(C_{\mathrm{corr}})$,
hence no ordering with respect to $P=I$ (with randomly rotated features it may be worse than $P=I$). These
values are \emph{local}: they rest on the weighted average $\bar w$ of the curvatures at $\etastar$
(Lemma~\ref{lem:B1-asymptotic}), whereas the rigorous trajectory bound \eqref{eq:B1-kappa-traj-stat} keeps the
smallest per-sample curvature, $m(\theta)\to\min_iw_i$, and need not improve on $P=I$. On the instances of
Figure~B1 it does not (comment after Corollary~\ref{cor:B1-kappa-traj}): any predicted gain of
$\hat\Sigma+\lambda I$ over $P=I$ rests on the local constants alone, and their use as a rate over a finite run
is heuristic. On ill-conditioned logistic instances ($\kappa\gg n$) with incoherent leverage scores
($d_\infty(\lambda)\ll n$), preconditioning should thus bring the \emph{effective} (local) regime back to
$\kappa_{\mathrm{eff}}\lesssim\frac{9n}{16}$, where the memory term $\frac1{3n}$ is the bottleneck
(Remark~\ref{rem:B1-memory-bottleneck}). This does not extend to the block-separable quadratic instances of
(A4), on which each $f_i$ depends only on its own block of $d/n$ coordinates: there each $f_i''$ has rank
$d/n$, so that $\lmin(P^{-1/2}F''P^{-1/2})\le\tr(P^{-1/2}F''P^{-1/2})/d\le R_P^2/n$, i.e.\
$\kappaP\ge n>\frac{9n}{16}$ for every $P$, with equality for $P=F''$ ($R_P^2=n$, $\mu_P=1$): the parameter
term still binds.
\item \emph{The other ratios.} When $\frac{3}{16\kappaP}<\frac1{3n}$ (for logistic regression as soon as
$\kappa>\frac{9n}{16}$, by Proposition~\ref{prop:B1-nofreelunch}; on the block-separable instances of (A4) for
every $P$, by (P3)),
$r_{\mathrm{min}}=s+(1-s)\frac{1-1/(3n)}{1-3/(16\kappaP)}$ exactly, with $s:=\norm{\delta}_P^2/V_{t-1}^P$ the
parameter share of $V^P$: $r_{\mathrm{min}}$ always lies within $\frac1{3n}$ of $1$, tends to $1$ iff $V^P$ is
dominated by the parameter term, and the min step costs
$-n\log r_{\mathrm{min}}\approx(1-s)\big(\frac13-\frac{3n}{16\kappaP}\big)$ in the exponent per pass. If the memory
errors $z_i^{(t-1)}-\fip(\etastar)\approx H_i\delta'$ were built from a single past error $\delta'$ aligned with
$\delta$, one would get $r_{\mathrm{var}}\approx1-\muEff^P/\LEff^P$ (close to $1$ when the effective condition
number is large, hence smaller under preconditioning) and a negative correlation of $a_{t,i}$ and $b_{t,i}$,
i.e.\ $r_{\mathrm{sq}}\le\frac12$, since
$r_{\mathrm{sq}}=\frac12+\sum_ia_{t,i}^\top P^{-1}b_{t,i}\big/\sum_i\big(\norm{a_{t,i}}_{P^{-1}}^2+\norm{b_{t,i}}_{P^{-1}}^2\big)$.
Exactly, $r_{\mathrm{init}}=r_{\mathrm{coco}}(1)\,\muEff^P(\theta_0)/R_P^2$ by \eqref{eq:B1-rinit-identity}, which is
expected to increase under preconditioning.
\item \emph{Check of the transfer.} Running Euclidean SAGA on $(\tilde f_i)$ with the same seed must
reproduce $P$-SAGA exactly, up to rounding (Proposition~\ref{prop:B1-transfer}(iii)); this is a sanity
check of the implementation, and the six ratios computed either way must coincide.
\end{enumerate}

\paragraph{Outcome of the experiments.}
In the experiment of Figure~B1 of the report ($n=1000$, $d=50$, $\kappa=10^3$, Gaussian features with
covariance spectrum $j^{-\alpha}$, i.e.\ covariance condition number $d^\alpha$, for $\alpha\in\{1,2\}$; this
$\alpha$ is not the decrement $\frac{3\mu_P}{16R_P^2}$ of \eqref{ineq:P-min}; $5$ sampling seeds,
$\gamma=1/(4R_P^2)$ with $R_P^2$ the upper bound in \eqref{eq:B1-logistic-RP}), the contraction exponents per
pass of $\norm{\theta_t-\etastar}_P^2$ (least-squares slopes over passes 2--20, seed-averaged) for $P=I$,
$P_{\mathrm J}$, $\hat\Sigma+\lambda I$, $F''(\etastar)$ are $1.000$, $0.9414$, $1.114$, $1.106$ ($\alpha=1$) and
$0.6104$, $0.2687$, $1.257$, $1.242$ ($\alpha=2$). As (P1) predicts, the global constants get the ordering
wrong: $\kappaP=1000$, $1899$, $5083$, $3623$ and $1000$, $2852$, $3639$, $2089$, computed like the step with the
upper bound of \eqref{eq:B1-logistic-RP} (the exact values of \eqref{eq:B1-RP-muP} are within $5.2\%$: $1000$,
$1898$, $5057$, $3611$ and $1000$, $2841$, $3459$, $2051$; all $\ge\kappa$), so $P=I$ gets the largest certified
exponent ($0.1875$).

The local constants give the observed ordering (up to the near tie of the dense metrics):
$\lmax(M_P)/\lmin(M_P)=25.17$, $25.67$, $2.801$, $1$ and $49.5$, $116.2$, $8.42$, $1$, and the bound with
$\lmin(M_P)$ in place of $\mu_P$ (heuristic: Lemma~\ref{lem:B1-asymptotic} supports it only as
$\theta\to\etastar$) gives $0.3256$, $0.2924$, $0.3333$, $0.3333$ and $0.1912$, $0.07705$, $0.3333$, $0.3333$.
Corollary~\ref{cor:B1-kappa-traj} does not account for this ordering: its bound for $\hat\Sigma+\lambda I$
exceeds $\kappa$ over the whole fitting window and tends to a value above its bound for $P=I$ (comment after
the corollary). What
accounts for it is Lemma~\ref{lem:B1-asymptotic}, rigorous as $\theta\to\etastar$ and used heuristically as a
rate over passes 2--20: along every direction, $\muEff^P(\theta)$ tends to a value $\ge\lmin(M_P)$. For
$\hat\Sigma+\lambda I$, $\lmin(M_P)=0.117$ and $0.109$ is the smallest over directions of a weighted average of
$w_1,\dots,w_n,1$, not their minimum: it is $16$ and $18$ times $\min_iw_i$, so that $R_P^2/\lmin(M_P)=166$ and
$103$, against $R^2/\lmin(F''(\etastar))=576$ and $981$ for $P=I$.

The dense metrics saturate the memory term $\frac13$, as in (P3); their observed exponents are $3.3$--$3.8$
times larger. Along the measured trajectories, the exact one-step contraction of
Proposition~\ref{prop:B3-representation} of the B3 note, transferred by Proposition~\ref{prop:B1-transfer} and
evaluated at the instrumented states of passes 2--20, splits this factor as follows when the inequalities are
undone in order~A of Proposition~\ref{prop:B3-attribution} ((I4), (I2), (I3), (I1), then the min step (I5);
logarithms of mean exponents, all seeds). The memory coefficient, $1-\frac1n(1-\frac23r_{\mathrm{sq}}r_{\mathrm{var}})\approx1-\frac{0.65}n$
here instead of $1-\frac1{3n}$ (through $r_{\mathrm{sq}}\approx0.51$--$0.53$, i.e.\ the square-norm step (I1), since
$r_{\mathrm{var}}\approx1$), accounts for about $55\%$ ($\alpha=1$) and $52\%$ ($\alpha=2$) of the logarithm, a
factor $1.95$--$1.99$. The min step (I5), which replaces by the slower memory rate a weighted mean in which the
parameter term (median share $2$--$7\%$ of $V^P$) contracts $7$--$21$ times faster, accounts for about $36\%$
and $39$--$40\%$, a factor $1.55$--$1.70$. The remaining $8$--$9\%$, a factor $1.10$--$1.12$, separates the
measured slopes from the exact one-step contraction (mean exponents $1.01$ and $1.11$--$1.12$ per pass) and is
pathwise: by Jensen's inequality, $\Ei[\log(V_t^P/V_{t-1}^P)]\le\log(\Ei[V_t^P]/V_{t-1}^P)$. Relative to the exact
one-step contraction, the memory coefficient and (I5) thus share the gap about $60/40$ ($\alpha=1$) and $57/43$
($\alpha=2$). None of this is a statement about $\E V_t^P$: the identity of Step~2 gives
$\E\sum_i\norm{z_i^{(t)}-\fip(\etastar)}_{P^{-1}}^2\ge(1-\frac1n)^t\sum_i\norm{z_i^{(0)}-\fip(\etastar)}_{P^{-1}}^2>0$,
the contribution of the slots never refreshed (an event of probability $(1-\frac1n)^t\approx e^{-t/n}$ per slot;
in the five runs every slot has been refreshed after $5.7$ to $9.1$ passes), so that no bound of the form
$\E V_t^P\le C\rho^tV_0^P$ holds with $\rho<1-\frac1n$. Such a bound certifies at most
$n\log\frac n{n-1}\approx1.0005$ per pass, $3$ times the local bound, whereas along the measured trajectories the
memory error itself decays with exponents $1.12$--$1.25$ per pass, faster than its expectation can (the exact
one-step contraction is faster than $1-\frac1n$ at the typical states, while $\E V_t^P$ is carried by the rare
states with stale slots): a factor $1.10$--$1.26$ of the observed exponents is invisible in the expectation
$\E V_t^P$ that the analysis controls.

The Jacobi metric is the slowest, which (P3) leaves open and Proposition~\ref{prop:B1-diag} allows: here
$\operatorname{cond}(C_{\mathrm{corr}})=52$ ($\alpha=1$) and $2866$ ($\alpha=2$), and
$\lmax(M_{P_{\mathrm J}})/\lmin(M_{P_{\mathrm J}})$ is $9.2$ and $13.8$ times that of $\hat\Sigma+\lambda I$. As in
(P2), $r_{\mathrm{coco}}$ rises from $0.1289$ ($P=I$) to $0.3288$ and $0.3377$, and $r_{\mathrm{sc}}$ falls from
$0.4233$ to $0.01692$ and $0.02288$ ($\alpha=1$, geometric means over passes 2--20).

The heuristics of (P4) fail for the dense metrics. $r_{\mathrm{var}}$ \emph{increases} under preconditioning
(geometric means $0.9936$, $0.9956$, $0.9990$, $0.9990$ for $\alpha=1$ and $0.958$, $0.970$, $0.9991$, $0.9984$
for $\alpha=2$): $1-r_{\mathrm{var}}$ is $36$--$173$ times smaller than $\muEff^P/\LEff^P$ for the dense metrics
($(0.9$--$1.6)\cdot10^{-3}$ against $0.035$--$0.036$ for $\alpha=1$ and $0.14$--$0.16$ for $\alpha=2$; geometric
means), and
within a factor $3.5$ of it for $P=I$ and $P_{\mathrm J}$. $r_{\mathrm{sq}}$ exceeds $\frac12$ (positive
correlation of $a_{t,i}$ and $b_{t,i}$) in $86$--$93\%$ of the instrumented states for $\hat\Sigma+\lambda I$ and
$F''(\etastar)$ (geometric means $0.525$, $0.526$ and $0.507$, $0.509$). Both failures occur in the
memory-bottleneck regime, where the stored gradients were computed at past iterates much farther from
$\etastar$ than $\theta_{t-1}$, contrary to the single aligned past error assumed in (P4). $r_{\mathrm{init}}$
increases as predicted ($8.5\cdot10^{-4}$, $1.2\cdot10^{-3}$, $3.7\cdot10^{-3}$, $7.3\cdot10^{-3}$ for $\alpha=1$),
and $r_{\mathrm{min}}$, an exact function of the parameter share $s$ of $V^P$ here (median $s=0.64$, $0.50$,
$0.066$, $0.063$ for $\alpha=1$), shows that with global constants the min step costs $0.23$--$0.27$ per pass for
the memory-dominated dense metrics. (P5) was not run as a separate experiment; the instrumentation in the
$P$-metric is checked instead against an exact enumeration of the $n$ possible steps, for a dense and a diagonal
$P$ (\texttt{tests/test\_vrlab.py}).
